% !TeX root = ../../Book1.tex
% 纲要标题：## 第17章　群作用、齐性空间与商结构
% 纲要位置：几何学计划纲要.md，第 768 行。
% 正文状态：初稿；尚未正式审核。

% 纲要内容开始
% 
% > 本章目标与证明范围：先证明轨道和商流形，再处理非自由作用。球面、射影空间和 Grassmann 流形承担切空间及不变结构的实际计算；有限稳定子商保留到局部模型层次。
% 
% 纲要内容结束

\chapter{群作用、齐性空间与商结构}
\label{g1:ch:17}
\BookChapterNumber{g1:ch:17}

\begin{chapterabstract}
	Lie 群作用把群中的连续运动变成流形上的向量场，稳定子则记录哪些运动在指定点上无法看见。本章先由闭子群构造齐性空间，把轨道的切空间写成 Lie 代数之商，再证明自由固有作用的轨道空间具有唯一相容的光滑结构。紧群上的归一化积分将直接由左不变密度构造，用于平均度量和固定点线性化。球面、射影空间与 Grassmann 流形提供切表示、图坐标和度量的具体计算。最后以固有作用的切片定理描述非自由作用，说明有限稳定子产生的局部商图为什么还须保存群数据。
\end{chapterabstract}

\begin{learninggoals}
	\begin{itemize}
		\item 构造基本向量场，核对左作用下的括号符号，证明稳定子、轨道秩与齐性空间切表示的关系。
		\item 从横截面和固有性证明自由固有作用的局部乘积及商流形定理，理解商映射的光滑泛性质，并辨别条件缺失的反例。
		\item 构造紧 Lie 群的归一化左右不变积分，用平均得到不变度量，并以稳定子表示组织齐性空间上的不变张量。
		\item 计算球面、实复射影空间及 Grassmann 流形的切空间、维数和不变度量，明确实维数与复维数的区别。
		\item 证明固定点线性化与固有作用切片定理，辨别底层商空间、有限稳定子数据、反射和非有效局部图。
	\end{itemize}
\end{learninggoals}

把平面中的点绕原点旋转，同一个圆周上的点就成为同一类。把球面上所有方向一起旋转，又可以把一个点送到任何另一个点。这两种观察分别引出轨道空间与齐性空间：\(M/G\) 用一点代表一整条轨道，而 \(G/G_p\) 用群参数描述经过 \(p\) 的那一条轨道。两个商的角色不同，前者保留轨道之间的方向，后者保留轨道自身的方向。

\ProseParagraphBreak

集合上能够分类，尚不足以得到光滑坐标。轨道可能稠密，稳定子可能改变大小，商中不同点也可能无法用不交邻域分开。本章的证明会分别说明自由性、固有性和稳定子紧性解决什么问题。特别是出现固定点时，商的底拓扑即使仍像一个圆盘，也未必保留全部几何数据。

\ProseParagraphBreak

本章的 Lie 群均为有限维实 Lie 群，光滑流形默认有限维、无边界、Hausdorff、第二可数；需要允许边界的商模型会另作说明。所有群作用默认左作用，基本向量场采用 \(\exp(tX)\cdot p\) 的正号初速度。复矩阵群及复参数空间先作为实流形使用，复维数在相应位置另外标明。

\ProseParagraphBreak

主要准备是\chapref{g1:ch:16}的指数映射、伴随表示与闭子群定理，\chapref{g1:ch:11}的逆函数及常秩定理，以及\chapref{g1:ch:12}的张量与度量。紧群积分使用\chapref{g1:ch:14}已经建立的密度积分和换元公式；这里会重新说明紧性怎样保证归一化与参数求导。切片证明采用紧稳定子的平均坐标。

\ProseParagraphBreak

延伸阅读可查 Lee 第二版第20章的 Infinitesimal Generators of Group Actions（第525--529页），以及第21章的 Quotients of Manifolds by Group Actions（第541--547页）和 Homogeneous Spaces（第550--554页）\cite[第525--529、541--547、550--554页]{Lee2012Smooth}。比较基本场公式时须先核对左、右作用及指数中的正负号；比较有限群商时则须核对文献是否允许反射和非有效局部群。

% !TeX root = ../../Book1.tex
% 纲要标题：### 17.1 作用和轨道的微分几何
% 纲要位置：几何学计划纲要.md，第 772 行。
% 正文状态：初稿；尚未正式审核。

% 纲要内容开始
% 
% * 无穷小作用与基本向量场
% * 稳定子、轨道切空间与秩
% * 齐性空间的光滑结构
% * 等变映射及其微分
% 
% 纲要内容结束

\section{作用和轨道的微分几何}
\label{g1:sec:1701}

在\secref{g1:sec:0104}中，群作用把集合分解为轨道，轨道又与稳定子群的陪集空间对应。现在群和被作用的空间都是流形，我们要把这份集合论描述变成微分几何：群中的一个切向量怎样推动点运动，轨道有多少个独立方向，陪集空间又怎样取得光滑坐标？这里必须保留一个区别：轨道可以有自己的流形拓扑，而这个拓扑未必就是环境流形给它的子空间拓扑。

\subsection{无穷小作用与基本向量场}
\label{g1:subsec:1701:01}

设 \(G\) 为 Lie 群，\(M\) 为光滑流形。\term[guanghuaqunzuoyong]{光滑群作用}[smooth group action]是光滑映射 \(a:G\times M\to M\)，满足 \(e\cdot p=p\) 和 \(g\cdot(h\cdot p)=(gh)\cdot p\)。本章默认左作用，记 \(a_g(p)=g\cdot p\)；\(a_g\) 是微分同胚，逆为 \(a_{g^{-1}}\)。右作用也可用同样方法处理，但变换次序与某些括号符号需要随之调整。

\ProseParagraphBreak

记 \(\mathfrak g=\mathrm{T}_eG\)。沿群中的单参数子群 \(\exp(tX)\) 运动，会使每个点 \(p\) 描出一条曲线。把其初速度收集起来，便得到 \(M\) 上的向量场。

\begin{definition}\label{g1:def:fundamental-vector-field}
	设 \(G\) 为有限维实 Lie 群，\(M\) 为光滑流形，\(a:G\times M\to M\) 为光滑左作用，记 \(g\cdot p=a(g,p)\)，\(\mathfrak g=\operatorname{Lie}(G)\)。对 \(X\in\mathfrak g\)，定义对应的\term[jibenxiangliangchang]{基本向量场}[fundamental vector field]为
	\begin{equation}\label{g1:eq:fundamental-field-convention}
		X_M(p)=\left.\frac{\mathrm{d}}{\mathrm{d}t}\right|_{t=0}\exp(tX)\cdot p.
	\end{equation}
	映射 \(X\mapsto X_M\) 称为该作用的\term[wuqiongxiaozuoyong]{无穷小作用}[infinitesimal action]，本章始终采用式中的正号约定。\symidx{\(X_M\)：正号约定下的基本向量场}
\end{definition}

\begin{proposition}[基本场的流与括号]\label{g1:prop:fundamental-fields}
	设有限维实 Lie 群 \(G\) 光滑左作用于无边界光滑流形 \(M\)。基本场 \(X_M\) 光滑且完备，其流为 \(\Phi_t(p)=\exp(tX)\cdot p\)。映射 \(X\mapsto X_M\) 线性，并满足
	\begin{equation}\label{g1:eq:fundamental-field-pushforward-bracket}
		(a_g)_*X_M=(\operatorname{Ad}_gX)_M,
		\qquad [X_M,Y_M]=-[X,Y]_M.
	\end{equation}
\end{proposition}
\begin{proof}
	对固定 \(p\)，定义轨道映射 \(\alpha_p:G\to M\)，\(\alpha_p(g)=g\cdot p\)。由于 \((\mathrm{d}\exp)_0=I\)，有 \(X_M(p)=(\mathrm{d}\alpha_p)_eX\)。微分对向量线性，并随 \(p\) 光滑变化，所以基本场光滑，且对 \(X\) 线性。

	单参数子群的乘法给出 \(\Phi_{t+s}=\Phi_s\circ\Phi_t\)。在 \(s=0\) 求导得 \(\partial_t\Phi_t(p)=X_M(\Phi_t(p))\)，初值为 \(p\)。此式对一切实数 \(t\) 都有定义，故它是全局流。

	由\cref{g1:prop:exponential-naturality}，
	\[
		g\exp(tX)\cdot p=\exp(t\operatorname{Ad}_gX)\cdot(g\cdot p).
	\]
	在 \(t=0\) 求导便得推前公式。为核对括号，回忆\cref{g1:prop:flow-pullback-derivative}：\([X_M,Y_M]=\left.\frac{\mathrm{d}}{\mathrm{d}t}\right|_0\Phi_t^*Y_M\)，其中场的拉回用逆向流的微分。因此
	\[
		\Phi_t^*Y_M=(a_{\exp(-tX)})_*Y_M
		=(\operatorname{Ad}_{\exp(-tX)}Y)_M.
	\]
	再用 \(\operatorname{Ad}_{\exp(tX)}=e^{t\operatorname{ad}_X}\)，在零点求导得到 \(-[X,Y]_M\)。
\end{proof}

负号来自左作用与\eqref{g1:eq:fundamental-field-convention}的正号约定。若改用 \(X\mapsto-X_M\)，就得到保持括号的映射。具体地，\(G\) 左乘自身时，基本场在 \(g\) 处是曲线 \(\exp(tX)g\) 的速度，因而是右不变场；前章的左不变场则沿 \(g\exp(tX)\) 运动。

\begin{example}\label{g1:exm:linear-action-fundamental}
	\(\operatorname{GL}(n,\mathbb R)\) 自然作用于 \(\mathbb R^n\) 时，\(X_M(v)=Xv\)。按向量场的坐标括号，
	\[
		[X_M,Y_M](v)=YXv-XYv=-(XY-YX)v,
	\]
	直接验证了负号。对平面旋转，取 \(J=\begin{pmatrix}0&-1\\1&0\end{pmatrix}\)，得到 \(J_M=-y\partial_x+x\partial_y\)。原点处这个场为零，其余点沿以原点为圆心的圆周运动。场在某点为零，意味着这一个单参数运动固定该点，并不意味着整个群都固定该点。
\end{example}

\subsection{稳定子、轨道切空间与秩}
\label{g1:subsec:1701:02}

固定 \(p\in M\)，记 \(G_p=\{g:g\cdot p=p\}\)、\(G\cdot p=\{g\cdot p:g\in G\}\)。这仍是\cref{g1:prop:orbit-stabilizer}中的稳定子群与轨道。回顾第10章的自由性：若每个 \(G_p=\{e\}\)，则称该作用为自由作用；若 \(\bigcap_{p\in M}G_p=\{e\}\)，则称该作用为\term[youxiaozuoyong]{有效作用}[effective action]。前者逐点排除非平凡稳定子，后者只排除处处不起作用的元素，二者不能混同。

\begin{proposition}[轨道映射的秩]\label{g1:prop:orbit-map-rank}
	设有限维实 Lie 群 \(G\) 光滑左作用于无边界光滑流形 \(M\)，\(p\in M\)。稳定子 \(H=G_p\) 是闭的嵌入 Lie 子群，其 Lie 代数为
	\begin{equation}\label{g1:eq:stabilizer-lie-algebra}
		\mathfrak h=\ker(\mathrm{d}\alpha_p)_e
		=\{X\in\mathfrak g:X_M(p)=0\}.
	\end{equation}
	轨道映射 \(\alpha_p\) 的秩恒等于 \(\dim G-\dim H\)，其在单位元处的微分像为 \(\{X_M(p):X\in\mathfrak g\}\)。此外，\(G_{g\cdot p}=gG_pg^{-1}\)。
\end{proposition}
\begin{proof}
	\(H=\alpha_p^{-1}(\{p\})\)，单点集在 Hausdorff 流形中闭，所以 \(H\) 闭。\cref{g1:thm:closed-subgroup}赋予它嵌入 Lie 子群结构，并且说明 \(X\in\mathfrak h\) 当且仅当 \(\exp(tX)\in H\) 对所有 \(t\) 成立。这一条件显然推出 \(X_M(p)=0\)。反之，若 \(X_M(p)=0\)，常值曲线 \(p\) 和曲线 \(\exp(tX)\cdot p\) 都是过 \(p\) 的积分曲线。由\secref{g1:sec:1301}的积分曲线唯一性，它们在整个共同定义域 \(\mathbb R\) 上相等，故 \(\exp(tX)\in H\)。这证明核的描述。

	恒等式 \(\alpha_p\circ L_g=a_g\circ\alpha_p\) 在单位元的微分给出
	\[
		(\mathrm{d}\alpha_p)_g\circ(\mathrm{d}L_g)_e
		=(\mathrm{d}a_g)_p\circ(\mathrm{d}\alpha_p)_e.
	\]
	两端的平移微分均可逆，故秩处处相同，等于 \(\dim\mathfrak g-\dim\mathfrak h\)。最后，\(k\) 固定 \(g\cdot p\) 当且仅当 \(g^{-1}kg\) 固定 \(p\)，这给出稳定子共轭公式。
\end{proof}

这个命题给出了轨道映射的秩及其核。核中的方向来自稳定子，因此可以先除去稳定子，把轨道参数化为一个陪集空间。常秩定理控制每个群元素附近的像；群中远处的元素仍可能把轨道送回当前邻域，所以陪集空间的流形拓扑与轨道像的子空间拓扑需要分别考察。

\subsection{齐性空间的光滑结构}
\label{g1:subsec:1701:03}

设 \(H\subseteq G\) 为闭子群。记 \(G/H\) 为左陪集 \(gH\) 的集合，赋予商映射 \(q(g)=gH\) 所决定的商拓扑。这里把同一个右 \(H\) 轨道压成一点；如果 \(H\) 不正规，不能把 \(G/H\) 当作商群。

\begin{theorem}[闭子群的齐性空间]\label{g1:thm:homogeneous-space}
	设 \(G\) 为有限维实 Lie 群，\(H\subseteq G\) 为闭子群。\(G/H\) 具有唯一的无边界、Hausdorff、第二可数光滑流形结构，使 \(q:G\to G/H\) 为光滑淹没。这个结构具有原来的商拓扑，维数为 \(\dim G-\dim H\)，左作用 \(G\times G/H\to G/H\) 光滑。\(q\) 在每点附近有光滑局部截面；对任意流形 \(N\) 和集合映射 \(f:G/H\to N\)，
	\[
		f\text{ 光滑}\quad\Longleftrightarrow\quad f\circ q\text{ 光滑}.
	\]
\end{theorem}
\begin{proof}
	先处理拓扑。若 \(A\subseteq G\) 开，则 \(q^{-1}(q(A))=AH=\bigcup_{h\in H}Ah\) 开，所以 \(q\) 是开映射。不同陪集 \(xH,yH\) 满足 \(x^{-1}y\notin H\)。由 \(H\) 闭及乘法连续，可取 \(x,y\) 的开邻域 \(A,B\)，使 \(A^{-1}B\cap H=\varnothing\)。于是 \(q(A)\) 与 \(q(B)\) 不交，证明商空间 Hausdorff。\(G\) 的一个可数开基经开映射 \(q\) 取像，就是 \(G/H\) 的可数开基。

	记 \(\mathfrak h=\mathrm{T}_eH\)，取线性补空间 \(\mathfrak g=\mathfrak m\oplus\mathfrak h\)。考虑
	\[
		F:\mathfrak m\times H\longrightarrow G,
		\qquad F(X,h)=\exp X\,h.
	\]
	它在 \((0,e)\) 的微分是 \((X,Y)\mapsto X+Y\)，为同构。逆函数定理给出 \(0,e\) 的邻域 \(W_0\subseteq\mathfrak m\)、\(V\subseteq H\)，使 \(F\) 在 \(W_0\times V\) 上是到开集的微分同胚。取更小的 \(W\subseteq W_0\)，使 \((\mathrm{d}F)_{(X,e)}\) 对 \(X\in W\) 都可逆，并且
	\[
		X,Y\in W,\quad \exp(-Y)\exp X\in H
		\quad\Longrightarrow\quad \exp(-Y)\exp X\in V.
	\]
	这一缩小可行，因为 \(H\) 取子空间拓扑，存在 \(G\) 中单位元邻域 \(D\) 满足 \(H\cap D\subseteq V\)，而 \(\exp(-Y)\exp X\) 在 \(X,Y\to0\) 时趋向 \(e\)。

	我们断言 \(F:W\times H\to (\exp W)H\) 是微分同胚。若 \(\exp X\,h=\exp Y\,k\)，则 \(u=\exp(-Y)\exp X=kh^{-1}\in V\)，从而 \(F(X,e)=F(Y,u)\)。局部单射性给出 \(X=Y\)、\(u=e\)，继而 \(h=k\)。故 \(F\) 单射。它在 \((X,e)\) 的微分可逆，利用第二因子和目标上的右平移，便知在 \((X,h)\) 处也可逆。因此它处处为局部微分同胚，像 \((\exp W)H\) 开；结合单射性，得到断言。

	现在 \(X\mapsto q(\exp X)\) 是 \(W\) 到开集 \(q((\exp W)H)\) 的同胚：在上述乘积坐标中，\(q\) 正是忘掉 \(H\) 因子的投影，故其逆坐标连续。把这张图用左平移移到任意 \(gH\) 附近，得到局部参数 \(X\mapsto q(g\exp X)\)。在两张图的交叠处，坐标转换由
	\[
		X\longmapsto k^{-1}g\exp X
		\longmapsto F_k^{-1}(k^{-1}g\exp X)
		\longmapsto\text{第一分量}
	\]
	给出，其中 \(F_k\) 是第二张图所用的横截乘积。各箭头均光滑，因而得到光滑图册。它的维数为 \(\dim\mathfrak m\)，且 \(q\) 在乘积坐标下为投影，所以是淹没。

	局部截面为 \(q(g\exp X)\mapsto g\exp X\)，记为 \(\sigma\)。若 \(f\circ q\) 光滑，则 \(f=(f\circ q)\circ\sigma\) 在每个这样的邻域光滑；反向由复合直接得到。左作用在局部写成 \((g,x)\mapsto q(g\sigma(x))\)，故光滑。

	最后，若另一光滑结构也使 \(q\) 为淹没，则由\cref{g1:thm:immersion-submersion-normal-form}，\(q\) 对该结构同样有光滑局部截面，并且是开映射。其拓扑因此仍为商拓扑。用两种结构各自的局部截面，恒等映射在两个方向都可写成一个局部截面与另一结构下的 \(q\) 的复合，故双向光滑，证明唯一性。
\end{proof}

定理中的最后一个性质说明：若先在 \(G\) 上定义光滑映射 \(F\)，并且证明它在每个陪集上常值，那么下降后的映射不必重新逐张坐标检查。下面的交换图同时表示 \(F=f\circ q\) 与 \(f\) 的唯一性：
\[
	\begin{tikzcd}
		G \arrow[r,"F"] \arrow[d,"q"'] & N \\
		G/H \arrow[ur,dashed,"f"'] &
	\end{tikzcd}
\]
局部截面把 \(F\) 的光滑性传递给 \(f\)。轨道映射恰好在稳定子的陪集上常值，因而可以用这个下降性质把它分解为陪集投影与轨道参数化。

\begin{proposition}[轨道的浸入结构]\label{g1:prop:orbit-immersion}
	设有限维实 Lie 群 \(G\) 光滑左作用于无边界光滑流形 \(M\)，\(p\in M\)。映射
	\[
		\overline\alpha_p:G/G_p\longrightarrow M,
		\qquad gG_p\longmapsto g\cdot p
	\]
	是单射浸入，像为 \(G\cdot p\)。把其定义域的光滑结构传给轨道后，
	\begin{equation}\label{g1:eq:orbit-tangent-space}
		\mathrm{T}_p(G\cdot p)=\{X_M(p):X\in\mathfrak g\}
		\cong\mathfrak g/\mathfrak g_p.
	\end{equation}
\end{proposition}
\begin{proof}
	集合上的单射性和像的描述正是轨道与陪集的对应。\(\alpha_p=\overline\alpha_p\circ q\)，所以前一定理的下降性质使 \(\overline\alpha_p\) 光滑。\((\mathrm{d}q)_e\) 满射，其核为 \(\mathfrak g_p\)；\cref{g1:prop:orbit-map-rank}给出 \((\mathrm{d}\alpha_p)_e\) 的同一个核。因此诱导的 \((\mathrm{d}\overline\alpha_p)_{eG_p}\) 单射。左平移和作用微分把此结论传到每点，证明浸入，并得到切空间公式。
\end{proof}

浸入轨道可能具有不同于像的子空间拓扑。例如 \(\mathbb R\) 按 \(t\cdot(x,y)=(x+t,y+\lambda t)\) 作用于 \(\mathbb R^2/\mathbb Z^2\)，其中 \(\lambda\) 无理。稳定子为零，轨道是\secref{g1:sec:1603}讨论过的稠密无理斜率线：其自身拓扑是一条实直线的拓扑，而远处参数会反复进入像中的任一小邻域。下一节的固有性将控制这种返回现象。

\ProseParagraphBreak

若 \(M\ne\varnothing\) 且只有一个轨道，则称该作用为\term[chuandizuoyong]{传递作用}[transitive action]；带有传递光滑 \(G\) 作用的流形称为 \(G\) 的\term[qixingkongjian]{齐性空间}[homogeneous space]。此时 \(G/G_p\) 与 \(M\) 集合上一一对应。传递性还使轨道映射满射；结合常秩定理与维数比较，可以进一步证明这个对应是微分同胚。

\begin{proposition}[传递作用的模型]\label{g1:prop:transitive-action-model}
	设有限维实 Lie 群 \(G\) 在非空无边界光滑流形 \(M\) 上传递光滑左作用，\(p\in M\)，\(G_p=\{g\in G:g\cdot p=p\}\)。则映射 \(\overline\alpha_p:G/G_p\to M\)、\(\overline\alpha_p(gG_p)=g\cdot p\) 为微分同胚。
\end{proposition}
\begin{proof}
	先说明满射的常秩映射 \(F:P\to M^n\) 必有秩 \(n\)，其中 \(P,M\) 为本章约定的流形。若秩 \(r<n\)，常秩定理给出定义域和像空间的局部坐标，使每个局部像包含在一个 \(r\) 维坐标平面中。用第二可数性及局部紧性，可取可数个紧坐标小球 \(K_j\) 覆盖 \(P\)，每个 \(K_j\) 都落在一张这样的定义域图内。于是 \(F(K_j)\) 在 \(M\) 中紧且闭，并因包含在低维坐标平面中而内部为空。

	非空流形不能被这些闭且内部为空的集合覆盖。这里可在 \(M\) 的一张坐标图内直接证明：取一个有内部的紧 Euclid 球，在其内部选择避开 \(F(K_1)\) 的更小闭球；再在后者内部选择避开 \(F(K_2)\) 的闭球，依次进行，并令半径趋于零。每一步可行，因为这些闭集内部为空。嵌套紧球有公共点，该点避开所有 \(F(K_j)\)，与覆盖矛盾。因此 \(F\) 的秩必须为 \(n\)。

	传递性使 \(\alpha_p:G\to M\) 满射，前面已证明其常秩，故它是淹没。于是 \(\dim(G/G_p)=\dim M\)。\(\overline\alpha_p\) 是同维流形之间的双射浸入，微分处处可逆，逆函数定理使其逆在每点附近光滑，从而它是微分同胚。
\end{proof}

为了以后计算不变度量，还需要知道稳定子怎样作用在一个切空间上。若 \(h\in G_p\)，则 \(a_h(p)=p\)，故 \((\mathrm{d}a_h)_p\) 是 \(\mathrm{T}_pM\) 的线性自同构。这给出\term[wendingzibiaoshi]{稳定子表示}[isotropy representation]
\[
	\rho_p:G_p\longrightarrow\operatorname{GL}(\mathrm{T}_pM),\qquad
	\rho_p(h)=(\mathrm{d}a_h)_p.
\]
链式法则保证 \(\rho_p(hk)=\rho_p(h)\rho_p(k)\)。它也保持轨道切空间，因此在法商空间 \(\mathrm{T}_pM/\mathrm{T}_p(G\cdot p)\) 上诱导线性作用。

\begin{proposition}[齐性空间的切表示]\label{g1:prop:homogeneous-tangent-isotropy}
	设 \(G\) 为有限维实 Lie 群，\(H\subseteq G\) 为闭子群，\(o=eH\in G/H\)，\(\mathfrak h=\mathrm{T}_eH\)。商映射的微分诱导同构 \(\mathrm{T}_o(G/H)\cong\mathfrak g/\mathfrak h\)。在此同构下，\(h\in H\) 的稳定子表示为
	\[
		\rho_o(h)(X+\mathfrak h)=\operatorname{Ad}_hX+\mathfrak h.
	\]
\end{proposition}
\begin{proof}
	\(q\) 是淹没，且局部乘积表明 \(\ker(\mathrm{d}q)_e=\mathfrak h\)，所以先得切空间同构。共轭保持 \(H\)，故 \(\operatorname{Ad}_h\) 保持 \(\mathfrak h\)，所写商空间算子有定义。再由
	\[
		h\exp(tX)H=\exp(t\operatorname{Ad}_hX)H
	\]
	在 \(t=0\) 求导，即得表示公式。
\end{proof}

\subsection{等变映射及其微分}
\label{g1:subsec:1701:04}

两个带有群作用的空间之间，普通光滑映射可以任意改变轨道。若希望映射保留对称性，就需要它与作用相容。设 \(G\) 作用于 \(M\)、\(K\) 作用于 \(N\)，\(\varphi:G\to K\) 为光滑同态。若光滑映射 \(F:M\to N\) 满足
\[
	F(g\cdot p)=\varphi(g)\cdot F(p).
\]
则称它为关于 \(\varphi\) 的\term[dengbianyingshe]{等变映射}[equivariant map]。若 \(G=K\)、\(\varphi=\operatorname{id}\)，就简称 \(G\) 等变。等变映射由一点的取值控制整条轨道，因此稳定子也会对这个取值施加限制。

\begin{proposition}[等变映射的微分]\label{g1:prop:equivariant-map-differential}
	设有限维实 Lie 群 \(G,K\) 分别光滑左作用于无边界光滑流形 \(M,N\)，作用映射记为 \(a_g(p)=g\cdot p\)、\(b_k(q)=k\cdot q\)。设 \(\varphi:G\to K\) 为光滑同态，\(F:M\to N\) 为满足 \(F(g\cdot p)=\varphi(g)\cdot F(p)\) 的光滑映射。对 \(p\in M\)，记两作用的稳定子为 \(G_p,K_{F(p)}\)，切表示为 \(\rho_p(h)=(\mathrm{d}a_h)_p\)、\(\rho_{F(p)}(k)=(\mathrm{d}b_k)_{F(p)}\)。则 \(\varphi(G_p)\subseteq K_{F(p)}\)，且对 \(X\in\operatorname{Lie}(G)\) 有
	\begin{equation}\label{g1:eq:equivariant-infinitesimal}
		(\mathrm{d}F)_pX_M(p)=((\mathrm{d}\varphi)_eX)_N(F(p)).
	\end{equation}
	若 \(h\in G_p\)，则其切表示满足
	\[
		(\mathrm{d}F)_p\circ\rho_p(h)
		=\rho_{F(p)}(\varphi(h))\circ(\mathrm{d}F)_p.
	\]
\end{proposition}
\begin{proof}
	若 \(h\cdot p=p\)，等变式给出 \(\varphi(h)\cdot F(p)=F(p)\)，故有稳定子包含。把 \(g=\exp(tX)\) 代入，并用指数自然性 \(\varphi(\exp(tX))=\exp(t(\mathrm{d}\varphi)_eX)\)，对 \(t\) 求导得\eqref{g1:eq:equivariant-infinitesimal}。最后对 \(F\circ a_h=a_{\varphi(h)}\circ F\) 在 \(p\) 求微分，得到切表示的相容式。
\end{proof}

在齐性空间上，等变映射由一个基点的像决定：其余点都由群作用得到。要使这个规定与基点的不同群参数表示相容，基点的稳定子必须固定所选像点。这样，构造映射的问题就转成了稳定子之间的包含关系。

\begin{proposition}[齐性空间之间的等变映射]\label{g1:prop:equivariant-homogeneous-maps}
	设 \(G\) 为有限维实 Lie 群，\(H,K\subseteq G\) 为闭子群。\(G/H,G/K\) 取使自然商投影为光滑淹没的商流形结构，并取 \(G\) 的标准左作用。\(G\) 等变光滑映射 \(G/H\to G/K\) 与满足 \(a^{-1}Ha\subseteq K\) 的陪集 \(aK\) 一一对应，映射为
	\[
		F_a(gH)=gaK.
	\]
	特别地，若 \(H\subseteq K\)，则 \(F_e:gH\mapsto gK\) 为光滑淹没，在基点的微分是自然满射 \(\mathfrak g/\mathfrak h\to\mathfrak g/\mathfrak k\)。
\end{proposition}
\begin{proof}
	若 \(F\) 等变且 \(F(eH)=aK\)，则 \(F(gH)=gaK\)，所以取值已被确定。\(H\) 固定 \(eH\)，故必须固定 \(aK\)，即 \(a^{-1}Ha\subseteq K\)。反之，在此包含条件下，\(g'=gh\) 给出 \(g'aK=ga(a^{-1}ha)K=gaK\)，所以公式与代表元无关。\(g\mapsto gaK\) 光滑，在 \(H\) 陪集上常值，由\cref{g1:thm:homogeneous-space}下降为光滑映射；等变性从公式直接得到。两个这样的映射相同，当且仅当它们在 \(eH\) 处的取值 \(aK\) 相同。

	若 \(H\subseteq K\)，关系 \(q_K=F_e\circ q_H\) 在单位元求微分，得所写商向量空间满射。左作用将满射性传到各点，故 \(F_e\) 是淹没。
\end{proof}

上述微分有清楚的几何含义：从 \(G/H\) 再压到 \(G/K\)，就是把原来仍能看见的 \(\mathfrak k/\mathfrak h\) 方向也视为等价。它的核给出被新等价关系消去的方向；保留下来的方向才是目标空间的切向量。下一节将把这种局部投影结构推广到不必传递的群作用。

% !TeX root = ../../Book1.tex
% 纲要标题：### 17.2 自由固有作用的商定理
% 纲要位置：几何学计划纲要.md，第 779 行。
% 正文状态：初稿；尚未正式审核。

% 纲要内容开始
% 
% * 固有性的局部和序列判据
% * 局部横截面与局部乘积
% * 商的 Hausdorff 性和光滑图册
% * 商映射的泛性质
% * 缺少自由或固有条件的反例
% 
% 纲要内容结束

\section{自由固有作用的商定理}
\label{g1:sec:1702}

把同一轨道上的点合并以后，横向方向应当留下来，沿轨道运动的方向则应被忘掉。自由性保证各条轨道具有同一个维数，但它不能阻止一条轨道反复回到很小的区域。为了让横向坐标真正成为商空间的坐标，还需要控制这种返回。本节的固有性正是一个同时限制出发点、到达点和群参数的条件。

\subsection{固有性的局部和序列判据}
\label{g1:subsec:1702:01}

\begin{definition}[固有作用]\label{g1:def:proper-lie-group-action}
	设有限维实 Lie 群 \(G\) 光滑左作用于有限维无边界光滑流形 \(M\)，两者均为 Hausdorff、第二可数。若映射
	\[
		\Theta:G\times M\longrightarrow M\times M,\qquad
		\Theta(g,p)=(gp,p)
	\]
	为固有映射，即每个紧集的原像都紧，则称这个作用为\term[guyouzuoyong]{固有作用}[proper action]。
\end{definition}

这里采用的是\cref{g1:def:proper-map-manifolds}的固有映射概念。需要约束的映射是 \(\Theta\)，而非单独的作用映射 \((g,p)\mapsto gp\)：后者忘掉出发点，即使直线上的平移也不满足同样的紧集原像条件。对两个子集 \(A,B\subseteq M\)，记
\[
	T(A,B)=\{g\in G:gA\cap B\ne\varnothing\}.
\]
它收集所有能够把 \(A\) 中某点送入 \(B\) 的群元素，下面称为从 \(A\) 到 \(B\) 的搬运集。这个记号中的先后次序与 \(\Theta\) 的两个输出次序不同：\(\Theta^{-1}(B\times A)\) 投影到群因子，才得到 \(T(A,B)\)。

\begin{proposition}[固有性的判据]\label{g1:prop:proper-action-criteria}
	设有限维实 Lie 群 \(G\) 光滑作用于有限维无边界光滑流形 \(M\)，且两者均为 Hausdorff、第二可数。对 \(A,B\subseteq M\)，令 \(T(A,B)=\{g\in G:gA\cap B\ne\varnothing\}\)。下列条件等价。
	\begin{equivalences}
		\item 作用固有。
		\item 对任意紧集 \(K,L\subseteq M\)，\(T(K,L)\) 在 \(G\) 中紧。
		\item 对任意 \(p,q\in M\)，存在开邻域 \(U\ni p\)、\(V\ni q\)，使 \(T(U,V)\) 在 \(G\) 中的闭包紧。
		\item 若 \(p_j\to p\)、\(g_jp_j\to q\) 在 \(M\) 中收敛，则 \((g_j)\) 有收敛子列。
	\end{equivalences}
	第二个条件也可只对 \(K=L\) 检查。
\end{proposition}
\begin{proof}
	先说明本证明使用序列的依据。流形有可数局部基，紧子集中的任意序列都有收敛子列：若无无限重复值，紧性给出一个聚点，再沿其递减局部基逐次选出子列。反过来，在第二可数空间的子集中，若每个序列都有收敛子列，则该子集紧。否则一个没有有限子覆盖的开覆盖可缩成可数子覆盖 \(O_1,O_2,\ldots\)，逐次选取 \(x_j\notin O_1\cup\cdots\cup O_j\)。任何收敛子列的极限都落在某个 \(O_k\)，从而子列最终也落在 \(O_k\)，矛盾。这使以下序列论证适用于本命题，而无须把它误用到任意拓扑空间。

	I 蕴含 II，因为 \(\Theta^{-1}(L\times K)\) 紧，而它在 \(G\) 中的投影就是 \(T(K,L)\)。反过来，设 II 成立，取紧集 \(C\subseteq M\times M\)，其两个投影分别记为 \(L,K\)。那么
	\[
		\Theta^{-1}(C)\subseteq T(K,L)\times K.
	\]
	右侧紧；\(C\) 在 Hausdorff 空间中闭，故左侧是右侧的闭子集，也紧，证明 I。若仅已知每个 \(T(K,K)\) 紧，则 \(T(K,L)\subseteq T(K\cup L,K\cup L)\)。此外 \(T(K,L)\) 总是闭的：若 \(g_j\to g\) 且 \(g_jx_j\in L\)、\(x_j\in K\)，紧性允许取 \(x_j\to x\in K\) 的子列，连续性给出 \(gx\in L\)，从而 \(g\in T(K,L)\)。流形中闭性可由序列检验，所以紧集的这个闭子集仍紧。

	II 蕴含 III，只须在 \(p,q\) 的坐标图中各取一个闭包紧的小开球。若反过来已知 III，对紧集 \(K,L\)，用相应乘积邻域 \(U\times V\) 覆盖 \(K\times L\)，再取有限子覆盖 \(U_i\times V_i\)。每个 \(g\in T(K,L)\) 都有见证点 \(x\in K\)、\(gx\in L\)；这对点落在某个 \(U_i\times V_i\)，所以
	\[
		T(K,L)\subseteq\bigcup_i\overline{T(U_i,V_i)}.
	\]
	右侧是有限个紧集之并。刚才已证左侧闭，故它紧。这也解释了为什么局部条件要允许任意一对 \(p,q\)，而不只检查同一点附近。

	I 蕴含 IV：序列 \((g_jp_j,p_j)\) 连同其极限组成紧集，固有性使全部 \((g_j,p_j)\) 落在一个紧集内，再取收敛子列即可。反之，若 IV 成立，对任意紧集 \(C\subseteq M\times M\) 及任意序列 \((g_j,p_j)\in\Theta^{-1}(C)\)，先由 \(C\) 紧取子列，使 \(p_j\to p\)、\(g_jp_j\to q\)。再由 IV 取 \(g_j\to g\) 的子列，便有 \((g_j,p_j)\to(g,p)\)。由于 \(C\) 闭且 \(\Theta\) 连续，这个极限仍在 \(\Theta^{-1}(C)\) 中。因此原像序列紧，按开头的说明即为紧集。
\end{proof}

若 \(G\) 离散，紧子集只有有限个点，于是第二个判据成为“每个紧集只被有限个群元素搬回与自身相交的位置”。这恢复了\secref{g1:sec:1003}讨论离散等距群时采用的固有性，但此时仍允许有限而非平凡的稳定子。

\begin{proposition}[固有作用的稳定子与轨道]\label{g1:prop:proper-action-orbits}
	设有限维实 Lie 群 \(G\) 固有地光滑作用于有限维无边界光滑流形 \(M\)，两者均为 Hausdorff、第二可数。则每个稳定子 \(G_p\) 紧，每条轨道 \(Gp\) 为 \(M\) 中的闭嵌入子流形，而且轨道等价关系
	\[
		R=\{(q,p)\in M\times M:q\in Gp\}
	\]
	在 \(M\times M\) 中闭。
\end{proposition}
\begin{proof}
	稳定子恰为 \(T(\{p\},\{p\})\)，因此紧。若 \((q_j,p_j)\in R\) 收敛到 \((q,p)\)，选取 \(g_j\) 使 \(q_j=g_jp_j\)。由序列判据，取子列后 \(g_j\to g\)，于是 \(q=gp\)，所以 \((q,p)\in R\)。这证明 \(R\) 闭。固定第二个变量为 \(p\)，\(Gp\) 就是连续映射 \(q\mapsto(q,p)\) 下 \(R\) 的原像，因而闭。

	为说明嵌入性，对任意紧集 \(C\subseteq M\)，轨道映射的原像 \(\alpha_p^{-1}(C)=T(\{p\},C)\) 紧。令 \(q:G\to G/G_p\)，则 \(\overline\alpha_p^{-1}(C)=q(\alpha_p^{-1}(C))\) 也紧。因此\cref{g1:prop:orbit-immersion}的单射浸入 \(\overline\alpha_p:G/G_p\to M\) 固有，由\cref{g1:prop:proper-injective-immersion}，它是闭像嵌入。这也说明在固有作用下，轨道的自然流形拓扑就是环境子空间拓扑。
\end{proof}

\begin{proposition}[紧群作用]\label{g1:prop:compact-group-action-proper}
	设紧的有限维实 Lie 群 \(G\) 光滑作用于有限维、Hausdorff、第二可数的无边界光滑流形 \(M\)。则作用固有。
\end{proposition}
\begin{proof}
	对紧集 \(C\subseteq M\times M\)，记第二个投影像为 \(K\)。\(\Theta^{-1}(C)\) 是紧集 \(G\times K\) 的闭子集，所以紧。
\end{proof}

\subsection{局部横截面与局部乘积}
\label{g1:subsec:1702:02}

在自由作用下，\cref{g1:prop:orbit-map-rank}给出单射 \(\mathfrak g\to \mathrm{T}_pM\)、\(X\mapsto X_M(p)\)，其像是轨道方向。选一个线性补空间，再用流形坐标把它实现成经过 \(p\) 的小子流形 \(S\)。这份\term[jubuhengjiemian]{局部横截面}[local transversal]的任务是为附近每条轨道提供一个代表。逆函数定理首先只保证小群参数内的代表唯一；固有性还将排除其他群元素把横截面送回自身。

\begin{lemma}[自由固有作用的局部乘积]\label{g1:lem:free-proper-local-product}
	设有限维实 Lie 群 \(G\) 自由且固有地光滑作用于 \(m\) 维无边界光滑流形 \(M\)，两者均为 Hausdorff、第二可数。令 \(r=\dim G\)。对每个 \(p\in M\)，存在经过 \(p\) 的 \(m-r\) 维嵌入子流形 \(S\)，使
	\begin{equation}\label{g1:eq:free-action-local-product}
		\Phi:G\times S\longrightarrow GS,\qquad \Phi(g,s)=gs
	\end{equation}
	为微分同胚，且饱和集 \(GS=\{gs:g\in G,\ s\in S\}\) 在 \(M\) 中开。每条与 \(S\) 相交的轨道只与 \(S\) 相交一次。
\end{lemma}
\begin{proof}
	先取 \(\mathrm{T}_pM\) 的直和分解
	\[
		\mathrm{T}_pM=\{X_M(p):X\in\mathfrak g\}\oplus E.
	\]
	在以 \(p\) 为原点的坐标图中，对坐标再作可逆线性变换，就可令最后 \(m-r\) 个坐标方向对应 \(E\)。前 \(r\) 个坐标置零，得到经过 \(p\) 的嵌入子流形 \(S_0\)，且 \(\mathrm{T}_pS_0=E\)。作用映射限制为 \(\Phi:G\times S_0\to M\)，其在 \((e,p)\) 的微分是
	\[
		\mathrm{d}\Phi_{(e,p)}(X,v)=X_M(p)+v.
	\]
	自由性和上述直和使它为同构。由\cref{g1:thm:inverse-function}，可以选取单位元开邻域 \(U\subseteq G\) 与 \(p\) 在 \(S_0\) 中的开邻域 \(S_1\)，使 \(\Phi|_{U\times S_1}\) 为到 \(M\) 中一个开集的微分同胚。此时 \((e,s)\)、\(s\in S_1\)，也都具有可逆微分。

	下一步是控制 \(U\) 外的群元素。我们断言，存在 \(p\) 在 \(M\) 中的开邻域 \(V\)，使
	\begin{equation}\label{g1:eq:proper-action-no-return}
		gV\cap V\ne\varnothing\quad\Longrightarrow\quad g\in U.
	\end{equation}
	若否，在 \(p\) 的一列递减坐标邻域 \(V_j\) 中，可选 \(s_j,t_j\to p\) 及 \(g_j\notin U\)，使 \(g_js_j=t_j\)。固有性的序列判据给出收敛子列 \(g_j\to g\)。取极限得 \(gp=p\)，自由性于是迫使 \(g=e\)。由于 \(U\) 是 \(e\) 的开邻域，这与 \(g_j\notin U\) 矛盾，证明断言。

	取 \(S=S_1\cap V\)，必要时再缩小为其中经过 \(p\) 的坐标邻域。如果 \(gs=t\)，其中 \(s,t\in S\)，则\eqref{g1:eq:proper-action-no-return}说明 \(g\in U\)。现在
	\[
		\Phi(g,s)=t=\Phi(e,t),
	\]
	而 \(\Phi\) 在 \(U\times S_1\) 上单射，故 \(g=e\)、\(s=t\)。若一般地 \(gs=ht\)，则 \(h^{-1}gs=t\)，由同一结论得 \(g=h\)、\(s=t\)。因此整个 \(G\times S\) 上的 \(\Phi\) 单射，而不只是小群邻域上的单射。

	还须说明这个全局于群因子的映射具有光滑逆。固定 \(g\in G\)，由
	\[
		\Phi(gk,s)=g\Phi(k,s)
	\]
	可知 \(\mathrm{d}\Phi_{(g,s)}\) 是 \(\mathrm{d}\Phi_{(e,s)}\) 与两侧平移微分同构的复合，所以处处可逆。逆函数定理使 \(\Phi\) 为局部微分同胚，因而为开映射，其像 \(GS\) 开。它到像又是双射，各局部光滑逆在交叠处由唯一性相同，便拼成全局光滑逆。
\end{proof}

\cref{g1:fig:free-action-product}显示这里的两个参数：群参数沿轨道变化，\(S\) 上的参数在不同轨道之间变化。图中画出的只是 \(G\) 的一小块；定理中的乘积却包含整个 \(G\)。因此它并不声称轨道本身像直线，也不要求 \(G\) 连通或紧。

\begin{figure}[htbp]
	\centering
	% !TeX root = ../Book1.tex
\begin{tikzpicture}[>=stealth,font=\small]
	\fill[blue!4] (-2,-1.2) rectangle (2,1.2);
	\foreach \y/\colour in {-0.8/blue!70!black,0/red!70!black,0.8/green!50!black}{
		\draw[thick,\colour,->] (-1.8,\y)--(1.8,\y);
		\fill[\colour] (0,\y) circle (2pt);
		\fill[\colour] (4,\y) circle (2pt);
	}
	\draw[thick] (0,-1.1)--(0,1.1) node[above] {\(\{e\}\times S\)};
	\draw[thick] (4,-1.1)--(4,1.1) node[above] {\(S\cong\pi(GS)\)};
	\draw[->] (2.3,0.2)--(3.4,0.2) node[midway,above] {\(\pi\)};
	\node at (0,-1.6) {\(G\times S\cong GS\)};
	\node at (4,-1.6) {横向坐标};
\end{tikzpicture}

	\caption[自由固有作用的局部乘积]{局部乘积 \(G\times S\cong GS\) 中的轨道方向与横向方向。每条水平线表示一条轨道的一部分，竖线表示横截面 \(\{e\}\times S\)；商映射忘掉群参数，只保留横向坐标。图中直线仅表示乘积坐标，不表示原流形中的轨道形状。}
	\label{g1:fig:free-action-product}
\end{figure}

\subsection{商的 Hausdorff 性和光滑图册}
\label{g1:subsec:1702:03}

给轨道集 \(M/G\) 赋予商拓扑，记 \(\pi(p)=[p]=Gp\)。对任意开集 \(O\subseteq M\)，有
\[
	\pi^{-1}(\pi(O))=GO=\bigcup_{g\in G}gO,
\]
它是开集，故 \(\pi\) 为开映射。这一点对任意连续群作用都成立；要得到流形，还须证明不同轨道能够分开，并为这些轨道配上相容的光滑坐标。

\begin{theorem}[自由固有作用的商流形]\label{g1:thm:free-proper-quotient}
	设 \(G\) 为 \(r\) 维实 Lie 群，\(M\) 为 \(m\) 维无边界光滑流形，两者均为 Hausdorff、第二可数。若 \(G\) 自由且固有地光滑作用于 \(M\)，则商拓扑下的 \(M/G\) 具有唯一的 \(m-r\) 维无边界光滑流形结构，使投影
	\[
		\pi:M\longrightarrow M/G
	\]
	成为光滑满射淹没。在\cref{g1:lem:free-proper-local-product}的每个局部乘积 \(G\times S\cong GS\) 中，\(\pi\) 对应第二因子投影，且 \(\pi|_S:S\to\pi(GS)\) 是微分同胚。
\end{theorem}
\begin{proof}
	先核验商空间的分离性与可数性。若 \([p]\ne[q]\)，则 \((q,p)\) 不在闭关系集 \(R\) 中。由\cref{g1:prop:proper-action-orbits}，可取开邻域 \(U\ni p\)、\(V\ni q\)，使 \((V\times U)\cap R=\varnothing\)。这就是说没有轨道同时经过 \(U\) 和 \(V\)，所以 \(\pi(U)\)、\(\pi(V)\) 是两个不交的开邻域，分别包含 \([p]\)、\([q]\)。故商空间 Hausdorff。取 \(M\) 的可数开基 \((B_j)\)，开集 \(\pi(B_j)\) 构成商空间的可数基：对 \([p]\in W\) 且 \(W\) 开，选 \(p\) 的基邻域 \(B_j\subseteq\pi^{-1}(W)\)，即可得到 \([p]\in\pi(B_j)\subseteq W\)。

	对局部横截面 \(S\)，记 \(W=\pi(GS)\)。\(GS\) 开且饱和，故 \(W\) 开；每条轨道与 \(S\) 交于唯一点，使 \(\pi|_S:S\to W\) 成为连续双射。它还开：若 \(A\subseteq S\) 开，则 \(G\times A\) 在 \(G\times S\) 中开，局部乘积微分同胚把它送到 \(GA\)，这是 \(M\) 的开集，所以 \(\pi(A)=\pi(GA)\) 在商中开。因此 \(\pi|_S\) 是同胚。取 \(S\) 上的光滑坐标 \(u:S\to\mathbb R^{m-r}\) 到一个开集，便得到 \(W\) 上的坐标 \(u\circ(\pi|_S)^{-1}\)。

	检验交叠时，设另一个横截面为 \(S'\)，局部乘积逆映射为 \(\Phi'^{-1}:GS'\to G\times S'\)。若 \(s\in S\cap GS'\)，唯一写成 \(s=g s'\)，其中 \(s'\in S'\)。用两张横截面表示同一轨道时，相应的转换正是
	\[
		S\cap GS'\longrightarrow S',\qquad
		s\longmapsto\operatorname{pr}_{S'}\bigl(\Phi'^{-1}(s)\bigr).
	\]
	定义域在 \(S\) 中开，右侧是光滑映射的复合，因此转换光滑；交换 \(S,S'\) 后得到光滑逆。于是上述坐标构成相容光滑图册，维数为 \(m-r\)，且没有边界。

	在乘积坐标中，\(\pi(gs)=\pi(s)\)，所以 \(\pi\) 就是 \((g,s)\mapsto s\)。它光滑且微分满射，因而为淹没；定义又使它满射。这些坐标还给出每个商点附近的光滑截面 \(\sigma:W\to M\)，即 \(\sigma([s])=s\)。

	最后证明唯一性。设同一商拓扑上还有另一份光滑结构，使同一个投影 \(\pi'\) 为光滑满射淹没。从刚构造的商到另一份商的恒等映射，在 \(W\) 上可写为 \(\pi'\circ\sigma\)，故光滑。反方向中，\cref{g1:thm:immersion-submersion-normal-form}说明 \(\pi'\) 的局部表达为投影；把被忘掉的坐标取常值，就得到目标点附近的光滑截面 \(\tau\)。从另一份商到本商的恒等映射在这片邻域上为 \(\pi\circ\tau\)，也光滑。两个恒等映射互逆，所以两份光滑结构相同。
\end{proof}

自由性使轨道方向始终具有维数 \(\dim G\)，固有性则既排除横截面的远处返回，也保证商的 Hausdorff 性。证明中第二可数性同样保留下来，不能只凭“每点附近看起来像 Euclid 空间”就称所得空间为本书意义的流形。

\begin{example}[正比例缩放的商]\label{g1:exm:positive-scaling-quotient}
	乘法 Lie 群 \((0,\infty)\) 作用于 \(\mathbb R^n\setminus\{0\}\)、\(n\ge1\)，作用为 \(a\cdot x=ax\)。它自由。若 \(x_j\to x\ne0\)、\(a_jx_j\to y\ne0\)，则
	\[
		a_j=\frac{\lVert a_jx_j\rVert}{\lVert x_j\rVert}
		\longrightarrow\frac{\lVert y\rVert}{\lVert x\rVert}>0,
	\]
	故作用固有。每条轨道恰与单位球面 \(S^{n-1}\) 相交一次，并且
	\[
		(0,\infty)\times S^{n-1}\longrightarrow\mathbb R^n\setminus\{0\},
		\qquad(a,u)\longmapsto au
	\]
	的光滑逆为 \(x\mapsto(\lVert x\rVert,x/\lVert x\rVert)\)。所以这里甚至有全局乘积，商就是 \(S^{n-1}\)，商投影对应 \(x\mapsto x/\lVert x\rVert\)。在 \(n=1\) 时，球面 \(S^0\) 为两个点；正负半轴分别是一条轨道，这也与零维商一致。
\end{example}

\subsection{商映射的泛性质}
\label{g1:subsec:1702:04}

拓扑商的泛性质见\cref{g1:prop:quotient-descent}：在每条轨道上取相同值的连续映射，才能成为商上的连续映射。现在局部横截面使同一结论提升到光滑范畴。实际使用时，这通常比直接计算商坐标更方便。

\begin{proposition}[光滑映射的下降]\label{g1:prop:quotient-smooth-universal-property}
	设有限维实 Lie 群 \(G\) 自由且固有地光滑作用于有限维无边界光滑流形 \(M\)，并设 \(N\) 为有限维无边界光滑流形；所有流形均为 Hausdorff、第二可数。给 \(M/G\) 以上一定理的光滑结构。则集合映射 \(h:M/G\to N\) 光滑，当且仅当 \(h\circ\pi:M\to N\) 光滑。

	特别地，光滑映射 \(f:M\to N\) 满足 \(f(gp)=f(p)\) 对所有 \(g,p\) 成立，当且仅当存在唯一光滑映射 \(\bar f:M/G\to N\)，使 \(f=\bar f\circ\pi\)。
\end{proposition}
\begin{proof}
	若 \(h\) 光滑，复合 \(h\circ\pi\) 光滑。反方向中，取商空间任意点附近的光滑截面 \(\sigma:W\to M\)，则
	\[
		h|_W=(h\circ\pi)\circ\sigma,
	\]
	所以它光滑。这样的 \(W\) 覆盖整个商，故 \(h\) 光滑；证明甚至不必预先假定 \(h\) 连续。

	若 \(f\) 在轨道上常值，定义 \(\bar f([p])=f(p)\)。轨道常值性保证定义不依赖代表元，满射性保证存在且唯一，刚证的光滑性判据保证 \(\bar f\) 光滑。反过来，任何复合 \(\bar f\circ\pi\) 在轨道上都取常值。
\end{proof}

这一泛性质可以画成
\[
	\begin{tikzcd}[column sep=large,row sep=large]
		M \arrow[r,"f\text{ 光滑且不变}"] \arrow[d,two heads,"\pi\text{ 淹没}"'] & N \\
		M/G \arrow[ur,dashed,"\bar f\text{ 唯一光滑}"'] &
	\end{tikzcd}
\]
图中的 \(N\) 是任意测试流形，不承担给定的 \(G\) 作用。若目标也带作用，等变映射未必沿源轨道取常值；要下降到轨道商，通常应同时把目标送到它的轨道商，不能直接套用这里的不变性条件。

\subsection{缺少自由或固有条件的反例}
\label{g1:subsec:1702:05}

\begin{example}[自由作用的非 Hausdorff 商]\label{g1:exm:irrational-flow-quotient}
	固定无理数 \(\alpha\)，令加法群 \(\mathbb R\) 作用于 \(T^2=\mathbb R^2/\mathbb Z^2\)：
	\[
		t\cdot[x,y]=[x+t,y+\alpha t].
	\]
	若某个 \(t\) 固定一点，则 \(t\in\mathbb Z\)、\(\alpha t\in\mathbb Z\)，所以 \(t=0\)，作用自由。\cref{g1:exm:irrational-torus-subgroup}已证明单位轨道是环面的真稠密子群；其他轨道是它的平移，也都稠密。取非空开集 \(O\subseteq T^2\)，每条轨道都与 \(O\) 相交，所以其饱和集 \(\mathbb R O\) 为整个环面。于是轨道商只有空集和全空间两个开集，而又有不止一条轨道，因此不是 Hausdorff 空间。

	这里固有性明确失败：紧集 \(K=T^2\) 的搬运集 \(T(K,K)=\mathbb R\) 非紧。每条轨道本身仍有一维浸入流形结构；失败的是把所有轨道合并以后，无法给它们彼此分离的邻域。
\end{example}

\begin{example}[旋转商的端点]\label{g1:exm:rotation-half-line-quotient}
	圆周群 \(S^1\) 通过复数乘法作用于 \(\mathbb C\cong\mathbb R^2\)。群紧，所以作用固有；原点的稳定子是整个 \(S^1\)，其他点的稳定子平凡，因此作用不自由。轨道由半径确定，\([z]\mapsto|z|\) 给出
	\[
		\mathbb C/S^1\cong[0,\infty).
	\]
	这是同胚：模长连续且在轨道上常值，故诱导连续映射；其逆 \(r\mapsto[r]\) 由正实半轴的连续包含给出。端点附近删去端点以后仍连通，而开区间中删去内点会出现两个分支，所以该商不是一维无边界流形。

	即使允许边界，商投影也不会满足前一定理的淹没结论。常用的光滑不变量 \(u(z)=|z|^2\) 在原点的微分为零。更一般地，若实值光滑函数 \(f\) 对旋转不变，则 \(f(-z)=f(z)\)，在原点求导得到 \(\mathrm{d}f_0=-\mathrm{d}f_0\)，因此 \(\mathrm{d}f_0=0\)。任何通过一维商坐标描述的光滑投影，在原点都不可能有秩 \(1\)。
\end{example}

非自由性也不意味着商一定出现坏点。例如 \(S^1\) 在任意流形 \(M\) 上的平凡作用是固有而不自由的，轨道商仍是 \(M\)，投影就是恒等映射。改变的是“轨道维数等于 \(\dim G\)”这一前提，故不能使用 \(\dim M-\dim G\) 的维数公式。下一步研究非自由作用时，需要保留稳定子及其横向作用，才能区分这种无害的冗余对称与旋转例中的端点。

% 纲要条目排为编号小节标题，并进入目录与书签；说明段仍保留于纲要注释中。

% !TeX root = ../../Book1.tex
% 纲要标题：### 17.3 不变结构与齐性模型
% 纲要位置：几何学计划纲要.md，第 787 行。
% 正文状态：初稿；尚未正式审核。

% 纲要内容开始
% 
% * 紧 Lie 群上的归一化不变积分
% 
%   以左不变密度及紧性构造积分，并用模函数的紧像证明右不变性；随即用于平均。本卷第19章进一步发展双不变度量，不反向承担此处引理。
% 
% * 紧群平均及不变度量
% * 稳定子表示与不变张量
% * 球面和射影空间的齐性表达
% * Grassmann 流形的切空间与维数
% 
% 纲要内容结束

\section{不变结构与齐性模型}
\label{g1:sec:1703}

球面的旋转既移动点，也保持各点的长度和夹角。一般群作用只规定怎样移动点，并不自动保持我们任意选取的度量。若群紧，可以把一个度量的所有变换平均起来；若作用传递，还可以只在一个点选取数据，再用群作用搬到全空间。这两种办法都需要解决选择无关性：平均要有不偏向某个群元素的积分，而沿群搬运要消除稳定子的影响。

\subsection{紧 Lie 群上的归一化不变积分}
\label{g1:subsec:1703:01}

有限群的平均是 \(|K|^{-1}\sum_{k\in K}\)。对圆周群，相应公式应为 \((2\pi)^{-1}\int_0^{2\pi}f(e^{it})\,\mathrm{d}t\)。两者都保持平移：重新选择群中的起点不改变平均值。一般紧 Lie 群没有一组整体坐标供我们直接积分，不过\cref{g1:prop:left-invariant-fields}已经说明，左平移能把单位元处的线性数据光滑地搬到全群；这次搬运的对象是一条正密度。

\begin{theorem}[紧群上的不变积分]\label{g1:thm:compact-group-integration}
	设 \(K\) 是有限维实紧 Lie 群，不要求连通。则 \(K\) 上存在唯一左不变的光滑正密度 \(\mu_K\)，满足 \(\int_K\mu_K=1\)。它还满足
	\[
		L_a^*\mu_K=\mu_K,\qquad R_a^*\mu_K=\mu_K,\qquad
		\iota^*\mu_K=\mu_K\quad(a\in K),
	\]
	其中 \(\iota(k)=k^{-1}\)。因此对任意连续实值函数 \(f\)，
	\begin{equation}\label{g1:eq:compact-group-integral-invariance}
		\begin{aligned}
			\int_K f(ak)\,\mathrm{d}\mu_K(k)
			&=\int_K f(ka)\,\mathrm{d}\mu_K(k)\\
			&=\int_K f(k^{-1})\,\mathrm{d}\mu_K(k)
			=\int_K f(k)\,\mathrm{d}\mu_K(k).
		\end{aligned}
	\end{equation}
	该积分线性、保持非负性，且 \(f\geq0\)、\(f\not\equiv0\) 时积分严格为正。
\end{theorem}
\begin{proof}
	先在 \(\mathfrak k=\mathrm{T}_eK\) 上任取一个正密度 \(\nu_e\)。它可以由任一正定内积的体积密度给出，见\cref{g1:prop:gram-volume-density}。对 \(v_1,\ldots,v_d\in \mathrm{T}_kK\) 定义
	\[
		\nu_k(v_1,\ldots,v_d)
		=\nu_e\bigl(\mathrm{d}(L_{k^{-1}})_kv_1,\ldots,
		\mathrm{d}(L_{k^{-1}})_kv_d\bigr),\qquad d=\dim K.
	\]
	左平移的微分及其逆随 \(k\) 光滑，所以 \(\nu\) 是光滑正密度。由 \(L_{(ak)^{-1}}\circ L_a=L_{k^{-1}}\)，有 \(L_a^*\nu=\nu\)。任意左不变密度都由单位元处的值唯一决定，而单位元处的密度空间是一维，故任意另一条左不变正密度都是 \(\nu\) 的正数倍。

	紧性使 \(\nu\) 本身具有紧支集。\cref{g1:thm:integral-well-defined}保证 \(c=\int_K\nu\) 有限且严格为正，故 \(\mu_K=\nu/c\) 是唯一满足归一化条件的左不变正密度。右不变性也可由左不变密度的唯一性得到：左、右平移交换，故 \(R_a^*\mu_K\) 也是左不变正密度。因而存在唯一正数 \(\Delta(a)\) 使
	\[
		R_a^*\mu_K=\Delta(a)\mu_K.
	\]
	在单位元的一组固定基上求值可知 \(\Delta\) 光滑。右平移复合及密度拉回的复合律给出 \(\Delta(ab)=\Delta(a)\Delta(b)\)。这是本处采用的\term[mohanshu]{模函数}[modular function]约定：记录右拉回对左不变密度的倍率。

	连续同态 \(\log\Delta:K\to(\mathbb R,+)\) 的像是紧加法子群。若其中有非零实数 \(r\)，它就含有所有整数倍 \(nr\)，与有界性矛盾。因此像只有零，\(\Delta\equiv1\)，得到右不变性。这一步同样适用于不连通的 \(K\)，没有把 \(K\) 换成单位元分支。

	由 \(\iota\circ L_a=R_{a^{-1}}\circ\iota\)，\(\iota^*\mu_K\) 是左不变正密度。\cref{g1:thm:oriented-change-variables}的密度换元公式表明 \(\int_K\iota^*\mu_K=\int_K\mu_K=1\)，故归一化的唯一性给出 \(\iota^*\mu_K=\mu_K\)。对三类微分同胚应用同一换元公式，就得到光滑 \(f\) 的\eqref{g1:eq:compact-group-integral-invariance}。

	连续 \(f\) 的积分也可直接用有限坐标覆盖和单位分解定义：逐片对连续紧支系数作 Euclid 积分。第14章证明中的坐标换元、共同细化与有限求和仍适用，所以积分和换元公式不变。非负性逐片成立；若 \(f(k_0)>0\)，连续性使它在一个坐标小邻域中有正下界，正密度便使这一块贡献正积分。零维时坐标积分就是点上的求和，论证也成立。
\end{proof}

以后写 \(\int_K f(k)\,\mathrm{d}\mu_K(k)\)，就是 \(\int_K f\mu_K\)，称为\term[guiyihuabubianjifen]{归一化不变积分}[normalized invariant integral]。这里积分的是密度 \(f\mu_K\)，无须选择定向。

\begin{example}\label{g1:exm:compact-integral-circle-finite}
	零维紧 Lie 群离散且紧，所以是有限群。左平移使各点的密度权重相同，归一化以后每点权重为 \(1/|K|\)，恰好恢复有限平均。对 \(S^1\)，角坐标中的密度 \((2\pi)^{-1}|\mathrm{d}t|\) 在换角度分支时相容，且被旋转保持；它的总积分为一，故就是上述 \(\mu_{S^1}\)。例如 \(m\in\mathbb Z\) 时
	\[
		\int_{S^1} z^m\,\mathrm{d}\mu_{S^1}(z)
		=\frac1{2\pi}\int_0^{2\pi}e^{imt}\,\mathrm{d}t
		=\begin{cases}1,&m=0,\\0,&m\ne0.\end{cases}
	\]
	复值积分按实部、虚部分别积分理解。紧性在一般定理中首先保证有限总量；在加法群 \(\mathbb R\) 上，左不变正密度必为 \(c|\mathrm{d}x|\)，其中 \(c>0\)，总积分无限，无法作同样的归一化。
\end{example}

\subsection{紧群平均及不变度量}
\label{g1:subsec:1703:02}

把光滑对象逐个平均，还须证明所得对象仍光滑。平均所用的群变量范围固定且紧，这正好提供了把参数微分移入积分所需的统一控制。以下引理是\subsecref{g1:subsec:1402:01}中紧区间参数积分论证的流形版本。

\begin{lemma}[紧参数积分]\label{g1:lem:compact-parameter-integral}
	设 \(K\) 是紧无边界光滑流形，\(\mu\) 为光滑密度，\(U\subset\mathbb R^m\) 开，\(F:K\times U\to\mathbb R\) 光滑。则
	\[
		I(x)=\int_K F(k,x)\,\mathrm{d}\mu(k)
	\]
	在 \(U\) 上光滑，且对每个多重指标 \(\alpha\)，有
	\[
		\partial_x^\alpha I(x)
		=\int_K\partial_x^\alpha F(k,x)\,\mathrm{d}\mu(k).
	\]
\end{lemma}
\begin{proof}
	给 \(K\) 取有限坐标覆盖和从属光滑单位分解。每一项都是对固定紧支集上的 Euclid 坐标变量积分，密度系数与单位分解均不依赖 \(x\)。固定 \(x_0\in U\)，取闭包紧且包含于 \(U\) 的小坐标盒 \(B\)。对第 \(j\) 个参数，微积分基本定理给出
	\[
		\frac{F(k,x+he_j)-F(k,x)}{h}
		=\int_0^1\partial_jF(k,x+the_j)\,\mathrm{d}t.
	\]
	当 \(x,x+he_j\in B\) 且 \(h\to0\) 时，右侧在紧参数集 \(K\times B\) 上局部一致趋于 \(\partial_jF(k,x)\)，因为该导数连续且在紧集上一致连续。各坐标片的密度系数绝对值可积，所以积分差的绝对值受“一致差上界乘有限积分”控制，差商可以移入积分。相同的一致控制还给出所得一阶导数的连续性。对任意 \(\partial_x^\alpha F\) 重复此论证，得到所有阶导数及其连续性。
\end{proof}

\begin{theorem}[平均度量]\label{g1:thm:compact-group-average}
	设紧 Lie 群 \(K\) 光滑地左作用于无边界光滑流形 \(M\)，记 \(a_k(p)=k\cdot p\)，\(\mu_K\) 为 \(K\) 上总质量为 \(1\) 的左右不变光滑正密度。给定 \(M\) 上任意光滑 Riemann 度量 \(g\)，公式
	\begin{equation}\label{g1:eq:averaged-metric}
		\overline g_p(v,w)
		=\int_K g_{k\cdot p}\bigl(\mathrm{d}(a_k)_pv,
		\mathrm{d}(a_k)_pw\bigr)\,\mathrm{d}\mu_K(k)
	\end{equation}
	定义一个 \(K\) 不变的光滑 Riemann 度量，即 \(a_h^*\overline g=\overline g\) 对每个 \(h\in K\) 成立。因此这样的作用总有不变 Riemann 度量，且不要求 \(M\) 紧。
\end{theorem}
\begin{proof}
	固定 \(p\) 时，各项都已拉回到 \(\mathrm{T}_pM\) 上的双线性型，故可以逐分量积分，而不是把不同切空间中的向量直接相加。双线性和对称性由积分的线性得到。若 \(v\ne0\)，微分同胚 \(a_k\) 的微分可逆，故 \(g_{k\cdot p}(\mathrm{d}a_kv,\mathrm{d}a_kv)>0\)。被积函数连续且处处正，\cref{g1:thm:compact-group-integration}给出 \(\overline g_p(v,v)>0\)。

	在 \(M\) 的一张坐标图内，把 \(v,w\) 依次取为坐标向量场。作用映射的微分随 \((k,p)\) 光滑，所得标量被积函数在 \(K\times U\) 上联合光滑。\cref{g1:lem:compact-parameter-integral}说明 \(\overline g\) 的每个局部分量光滑。最后，由 \(a_k\circ a_h=a_{kh}\) 得
	\[
		a_h^*\overline g
		=\int_K a_h^*a_k^*g\,\mathrm{d}\mu_K(k)
		=\int_K a_{kh}^*g\,\mathrm{d}\mu_K(k)
		=\overline g,
	\]
	最后一步用右不变积分。任意初始 \(g\) 的存在由\cref{g1:thm:bundle-metric-existence}用于 \(\mathrm{T}M\) 保证。
\end{proof}

同样可以平均协变张量、函数和微分形式，因为它们拉回后都属于固定纤维中的向量空间。对一般张量，平均未必保留非线性约束：例如一个非零向量场的平均可能为零，一个非退化二形式的平均也可能退化。度量的正定性能够保留，原因是正定双线性型构成凸集。

\begin{example}\label{g1:exm:rotation-average-metric}
	在 \(\mathbb R^2\) 上取常系数度量 \(g=a\,\mathrm{d}x^2+2b\,\mathrm{d}x\,\mathrm{d}y+c\,\mathrm{d}y^2\)，其中 \(a>0\)、\(ac-b^2>0\)。记其矩阵为 \(Q\)，平面旋转群的平均矩阵为
	\[
		\frac1{2\pi}\int_0^{2\pi}R_t^{\mathsf T}QR_t\,\mathrm{d}t
		=\frac{a+c}{2}I.
	\]
	这是因为 \(\cos^2t\)、\(\sin^2t\) 的平均均为 \(1/2\)，\(\sin t\cos t\) 的平均为零。平均消除了原来的优先方向，但保留了正定性。这里的底空间 \(\mathbb R^2\) 非紧，并不妨碍对紧的旋转参数积分。
\end{example}

\subsection{稳定子表示与不变张量}
\label{g1:subsec:1703:03}

在齐性空间 \(G/H\) 上，任何点都可写成 \(gH\)。若在基点 \(o=eH\) 选定一个内积，似乎只要用 \(g\) 把它搬到 \(gH\) 即可；但是同一点也等于 \(ghH\)。这两次搬运相差 \(h\) 在 \(\mathrm{T}_o(G/H)\) 上的作用，因此内积必须被整个稳定子保持。\cref{g1:prop:homogeneous-tangent-isotropy}已把这一作用写成
\[
	\mathrm{T}_o(G/H)\simeq\mathfrak g/\mathfrak h,\qquad
	\rho(h)(X+\mathfrak h)=\operatorname{Ad}_hX+\mathfrak h.
\]
此处 \(H\) 是闭 Lie 子群，可以不连通；只检查 \(\mathfrak h\) 的无穷小作用未必控制它的其余分支。

\begin{proposition}[不变张量的基点数据]\label{g1:prop:invariant-tensor-isotropy}
	设 \(G\) 为有限维实 Lie 群，\(H\subset G\) 为闭子群，\(o=eH\)，并给 \(G/H\) 取\cref{g1:thm:homogeneous-space}的光滑结构。对任意非负整数 \(r,s\)，取基点值给出如下两个实向量空间的线性同构：\(G/H\) 上的光滑 \(G\) 不变 \((r,s)\) 型张量场，以及 \(T^r_s(\mathrm{T}_o(G/H))\) 中被稳定子表示 \(\rho(H)\) 固定的张量。

	特别地，\(G\) 不变 Riemann 度量恰好对应 \(\mathfrak g/\mathfrak h\) 上的 \(\operatorname{Ad}(H)\) 不变正定内积。
\end{proposition}
\begin{proof}
	记 \(\ell_g(xH)=gxH\)。若 \(A:V\to W\) 是线性同构，张量的搬运在向量因子上用 \(A\)，在对偶因子上用 \((A^{-1})^*\)；记这个同构为 \(A_\#\)。它保持张量积、收缩和复合。对协变二阶张量，公式具体为 \((A_\#b)(v,w)=b(A^{-1}v,A^{-1}w)\)。

	若 \(T\) 是不变张量场，则 \(T_{gH}=(\mathrm{d}\ell_g)_\#T_o\)，所以它由 \(T_o\) 唯一决定。令 \(g=h\in H\)，因 \(hH=o\)，便知 \(T_o\) 被 \(\rho(h)=\mathrm{d}(\ell_h)_o\) 固定。反过来，给定这样的固定张量 \(t\)，规定
	\[
		T_{gH}=(\mathrm{d}(\ell_g)_o)_\#t.
	\]
	若改用代表 \(gh\)，则 \(\ell_{gh}=\ell_g\circ\ell_h\)，所得值为 \((\mathrm{d}\ell_g)_\#\rho(h)_\#t=(\mathrm{d}\ell_g)_\#t\)，因此良定义。

	还需证明它光滑。商投影 \(q:G\to G/H\) 在每点附近有光滑局部截面 \(s:U\to G\)，这是\cref{g1:thm:homogeneous-space}的局部乘积坐标所给出的结论。在 \(U\) 上，刚才的定义写成 \(T_x=(\mathrm{d}(\ell_{s(x)})_o)_\#t\)。作用映射及其微分光滑，逆微分也光滑，所以这些局部分量光滑。群乘法直接给出 \((\mathrm{d}\ell_a)_\#T_x=T_{ax}\)，即不变性。两种构造互逆，且线性。最后，同构搬运保持对称性和正定性，故正定内积恰好给出不变 Riemann 度量。
\end{proof}

\begin{corollary}[紧稳定子的度量]\label{g1:cor:compact-isotropy-metric}
	设 \(G\) 为有限维实 Lie 群，\(H\subset G\) 为紧子群，则 \(G/H\) 存在光滑 \(G\) 不变 Riemann 度量。这里不要求 \(G\) 紧。
\end{corollary}
\begin{proof}
	紧子群在 Hausdorff 的 \(G\) 中闭，故商流形和稳定子表示都有定义。在 \(\mathfrak g/\mathfrak h\) 上任取正定内积 \(b\)，令
	\[
		\overline b(u,v)
		=\int_H b\bigl(\rho(h)u,\rho(h)v\bigr)\,\mathrm{d}\mu_H(h).
	\]
	与\cref{g1:thm:compact-group-average}相同，线性、正性和右不变积分保证它为 \(H\) 不变正定内积。前一命题把它唯一延拓到 \(G/H\)。
\end{proof}

紧性在这里是一个充分条件。比如 \(G=\mathbb R^n\) 平移作用于自身，稳定子只有单位元，任何常系数正定内积都给出不变度量，尽管 \(G\) 非紧。相反，\(\mathbb R_{>0}\) 以伸缩作用于 \(\mathbb R\) 时不可能有不变 Riemann 度量：原点被固定，若其内积值为 \(b>0\)，伸缩 \(x\mapsto2x\) 的不变性便要求 \(4b=b\)。这说明在固定点处，稳定子的线性作用已经可能阻止不变度量的存在。

\subsection{球面和射影空间的齐性表达}
\label{g1:subsec:1703:04}

以下例子把商空间写法变成切空间和度量的具体计算。\cref{g1:prop:matrix-classical-lie-groups}已经构造了正交群、特殊正交群和酉群的 Lie 群结构。它们还是紧群：矩阵方程定义闭集，而每一列长度为一使所有分量有界。稳定子的块对角写法则来自“保持指定向量”或“保持指定子空间”这两个不同条件。

\begin{proposition}[球面模型]\label{g1:prop:sphere-homogeneous-model}
	对 \(n\geq1\)，映射
	\[
		\operatorname{SO}(n+1)/\operatorname{SO}(n)\longrightarrow S^n,
		\qquad AH\longmapsto Ae_0,
	\]
	是与 \(\operatorname{SO}(n+1)\) 作用相容的微分同胚。此处 \(\operatorname{SO}(n)\) 嵌入为 \(\operatorname{diag}(1,B)\)。在基点 \(e_0\)，商切空间通过
	\[
		\begin{pmatrix}0&-v^{\mathsf T}\\v&B\end{pmatrix}
		+\mathfrak{so}(n)\longmapsto v\in\mathbb R^n
	\]
	识别为 \(\mathrm{T}_{e_0}S^n=e_0^\perp\)，稳定子作用为 \(v\mapsto Bv\)。其标准内积延拓得到单位球面的诱导度量。
\end{proposition}
\begin{proof}
	任意单位向量 \(p\) 都能补成正交单位基。因 \(n\geq1\)，必要时反转另一列，可把该基的行列式调为 \(1\) 而保持第一列为 \(p\)，故作用传递。固定 \(e_0\) 的正交矩阵必保持 \(e_0^\perp\)，因此恰为 \(\operatorname{diag}(1,B)\)、\(B\in\operatorname{SO}(n)\)。矩阵作用在球面坐标中光滑，由\cref{g1:prop:transitive-action-model}得到所述微分同胚。

	轨道映射 \(A\mapsto Ae_0\) 的单位元微分是 \(X\mapsto Xe_0\)。反对称矩阵按上述块式写出后，微分取出 \(v\)，核正是右下角的 \(\mathfrak{so}(n)\)。稳定子在 \(e_0^\perp\) 上的作用是其右下角矩阵，因而保持标准内积。所有 \(A\in\operatorname{SO}(n+1)\) 都保持环境 Euclid 内积，所以球面诱导度量正是由基点内积延拓得到的那个不变度量。
\end{proof}

当 \(n=1\) 时，\(\operatorname{SO}(1)=\{1\}\)，上式就是 \(S^1\simeq\operatorname{SO}(2)\)。当 \(n=0\) 时，\(S^0\) 有两个点，\(\operatorname{SO}(1)\) 不能传递作用，不能原样套用该式；改用 \(\operatorname{O}(1)\) 便有 \(S^0\simeq\operatorname{O}(1)/\{1\}\)。

\ProseParagraphBreak

实射影空间的光滑图已在\cref{g1:exm:projective-smooth-atlas}中建立。复射影空间 \(\mathbb{CP}^n\) 同样由 \(\mathbb C^{n+1}\) 中的一维复子空间组成；这里把它视为实流形。给它单位球面按复相位 \(z\sim\lambda z\)、\(|\lambda|=1\) 取商的拓扑，在 \(z_i\ne0\) 处取坐标 \((z_j/z_i)_{j\ne i}\in\mathbb C^n\)。坐标逆把 \(1\) 插入第 \(i\) 位，再归一化到球面并取商，因而连续。换图公式仍是各分量除以一个非零分量，是实光滑映射；\(\mathbb C^n\simeq\mathbb R^{2n}\) 给出实维数 \(2n\)。映射 \([z]\mapsto zz^*/(z^*z)\) 是到 Hermitian 矩阵空间的连续单射，不同投影可用不交开球分离，故商 Hausdorff；有限张坐标图保证第二可数。于是这些图确实给出一个紧光滑流形。

\begin{proposition}[射影空间模型]\label{g1:prop:projective-homogeneous-models}
	对 \(n\geq0\)，有等变微分同胚
	\begin{align*}
		\mathbb{RP}^n&\simeq\operatorname{O}(n+1)/(\operatorname{O}(1)\times\operatorname{O}(n)),\\
		\mathbb{CP}^n&\simeq\operatorname{U}(n+1)/(\operatorname{U}(1)\times\operatorname{U}(n)).
	\end{align*}
	分母按相应的块对角矩阵嵌入，\(\operatorname{O}(0)\)、\(\operatorname{U}(0)\) 均理解为平凡群。统一记 \(\mathbb F=\mathbb R\) 或 \(\mathbb C\)。在一条直线 \(L\) 处，有自然识别
	\[
		\mathrm{T}_L\mathbb{FP}^n\simeq\operatorname{Hom}_{\mathbb F}(L,\mathbb F^{n+1}/L)
		\simeq\operatorname{Hom}_{\mathbb F}(L,L^\perp),
	\]
	第二个识别使用标准内积。对 \(L=\mathbb Fe_0\)，将 \(A\) 识别为 \(v=Ae_0\)，稳定子 \(\operatorname{diag}(\lambda,B)\) 作用为 \(v\mapsto Bv\lambda^{-1}\)。实内积 \(\operatorname{Re}(v^*w)\) 因而给出射影空间上的不变 Riemann 度量。
\end{proposition}
\begin{proof}
	正交或酉变换可以把任意单位向量送到另一单位向量，所以在直线上传递。固定 \(\mathbb Fe_0\) 的变换也保持其正交补，恰为所列块对角矩阵。矩阵作用在射影图中由线性分量相除写出，分母在选定图域非零，故光滑。由\cref{g1:prop:transitive-action-model}，两种集合上的轨道对应都提升为所述微分同胚。

	先在 \(L\) 的某个补空间 \(Q\) 中描述邻近直线。凡投影到 \(L\) 非零的直线都唯一写成 \(\{w+A(w):w\in L\}\)，其中 \(A\in\operatorname{Hom}_{\mathbb F}(L,Q)\)。这正是以 \(L\) 为坐标原点的射影图。因此过 \(L\) 的曲线 \(A(t)\)、\(A(0)=0\) 的切向量为 \(A'(0)\)。把 \(A'(0)w\) 取模 \(L\)，便消除了补空间选择，得到第一个自然识别；取 \(Q=L^\perp\) 得第二个识别。

	稳定子将图中向量 \(w+A(w)\) 送到 \(\lambda w+B A(w)\)，改以 \(w'=\lambda w\) 作参数，新的线性映射为 \(B A\lambda^{-1}\)。于是基点切向量变为 \(Bv\lambda^{-1}\)。因为 \(B\) 正交或酉且 \(|\lambda|=1\)，它保持 \(\operatorname{Re}(v^*w)\)。\cref{g1:prop:invariant-tensor-isotropy}给出不变度量。
\end{proof}

这个度量还能直接算到坐标中。取射影图 \([1:z]\)，\(z\in\mathbb F^n\)，以 \(u,v\in\mathbb F^n\) 表示坐标切向量，令 \(r=1+z^*z\)。曲线 \([1:z+tu]\) 的代表向量为 \(w(t)=(1,z+tu)\)，速度为 \((0,u)\)。先把速度投影到 \(w(0)^\perp\)，再除以 \(|w(0)|\)，便得到定义在单位线向量上的切映射。故
\begin{equation}\label{g1:eq:projective-invariant-metric}
	g_z(u,v)=\operatorname{Re}\frac{r\,u^*v-(u^*z)(z^*v)}{r^2}.
\end{equation}
分母中的两次归一化分别来自正交投影和线向量的长度。由 Cauchy--Schwarz 不等式，分子在 \(u=v\ne0\) 时至少为 \(|u|^2>0\)，所以公式也直接显示正定性。实情形的 \(S^n\to\mathbb{RP}^n\) 是局部等距映射：单位代表 \(p\) 处的速度 \(u\perp p\) 对应 \(A(p)=u\)，内积没有改变。复情形还要去掉相位方向 \(ip\)，因此不能把奇维球面上的所有切向量直接当作射影切向量。

\subsection{Grassmann 流形的切空间与维数}
\label{g1:subsec:1703:05}

\cref{g1:def:grassmannian}已经把 \(\operatorname{Gr}(k,V)\) 定义为 \(V\) 中所有 \(k\) 维线性子空间的集合，\subsecref{g1:subsec:0206:02}还用矩阵的非零子式给出了局部参数。现在要说明这些参数怎样组成光滑流形，以及一个子空间的“无穷小移动”是什么。我们继续允许 \(\mathbb F=\mathbb R\) 或 \(\mathbb C\)，并把 \(k=0\) 的情形补入约定：它只有零子空间这一个点。

\begin{theorem}[Grassmann 流形]\label{g1:thm:grassmann-homogeneous-model}
	设 \(\mathbb F=\mathbb R\) 或 \(\mathbb C\)，\(n\ge0\) 及 \(0\leq k\leq n\) 均为整数，\(V=\mathbb F^n\)，分别赋予标准 Euclid 内积或标准 Hermitian 内积。则 \(\operatorname{Gr}(k,V)\) 有自然的紧光滑流形结构，并有等变微分同胚
	\begin{align*}
		\operatorname{Gr}(k,\mathbb R^n)&\simeq\operatorname{O}(n)/(\operatorname{O}(k)\times\operatorname{O}(n-k)),\\
		\operatorname{Gr}(k,\mathbb C^n)&\simeq\operatorname{U}(n)/(\operatorname{U}(k)\times\operatorname{U}(n-k)).
	\end{align*}
	它们的实维数分别为 \(k(n-k)\)、\(2k(n-k)\)。在任意 \(W\in\operatorname{Gr}(k,V)\) 处，有自然同构
	\begin{equation}\label{g1:eq:grassmann-tangent}
		\mathrm{T}_W\operatorname{Gr}(k,V)
		\simeq\operatorname{Hom}_{\mathbb F}(W,V/W).
	\end{equation}
	复情形右边的复向量空间在此先作为实向量空间理解；其复维数为 \(k(n-k)\)。
\end{theorem}
\begin{proof}
	记 \(E=\mathbb F^k\oplus0\)、\(Q=0\oplus\mathbb F^{n-k}\)。给任意 \(k\) 维子空间选正交单位基并补成 \(V\) 的正交单位基，可见正交群或酉群在所有这样的子空间上传递。保持 \(E\) 的等距线性映射也保持 \(Q=E^\perp\)，稳定子恰好为块对角群 \(H\)。它是闭子群，所以\cref{g1:thm:homogeneous-space}先在 \(G/H\) 上给出光滑结构。集合双射 \(gH\mapsto gE\) 将这个结构赋予 Grassmann 空间；\(G\) 紧，故它也紧。

	为了辨认这套结构，考虑所有投影到 \(E\) 为同构的子空间。它们恰好是线性映射 \(A:E\to Q\) 的图：
	\begin{equation}\label{g1:eq:grassmann-graph-chart}
		W_A=\{x+Ax:x\in E\}.
	\end{equation}
	存在性是先用到 \(E\) 的投影把 \(W\) 参数化，唯一性来自每个 \(x\) 只有一个相应图向量。若 \(g\in G\) 的前 \(k\) 列按 \(E\oplus Q\) 写成 \(\binom{a}{c}\)，则此投影可逆恰好等价于 \(a\) 可逆，此时 \(A=ca^{-1}\)。该条件在 \(G\) 中开，且对右乘 \(H\) 不变；商投影是开映射，所以这些子空间组成 \(G/H\) 的开集。用商的光滑局部截面代入 \(ca^{-1}\)，可知图坐标光滑。

	反向，对块矩阵 \(\left(\begin{smallmatrix}I&0\\A&I\end{smallmatrix}\right)\) 的列依次作正交归一化，得到 \(g(A)\in G\)，其前 \(k\) 列仍张成 \(W_A\)。这个矩阵处处可逆，故每一步减去前面投影后的向量不为零；各步只用光滑加减、内积与正数平方根，因此 \(g(A)\) 光滑依赖 \(A\)。于是 \(A\mapsto g(A)H\) 光滑，确为上述坐标的逆。更换 \(E\oplus Q\) 的坐标后，若同一个向量的分量按固定可逆块矩阵 \(\left(\begin{smallmatrix}a&b\\c&d\end{smallmatrix}\right)\) 变换，则
	\begin{equation}\label{g1:eq:grassmann-chart-transition}
		A'=(c+dA)(a+bA)^{-1}.
	\end{equation}
	交叠条件恰是 \(a+bA\) 可逆，因此转换光滑；复情形还由复有理函数给出。各坐标 \(k\) 平面对应的这类图覆盖全空间，因为任意秩为 \(k\) 的基矩阵总有一个非零 \(k\) 阶子式。这也恢复了第2章的矩阵参数，说明上述结构与那些图坐标相容。图空间 \(\operatorname{Hom}_{\mathbb F}(E,Q)\) 给出所述维数。

	最后在任意 \(W\) 处选补空间并取以 \(W\) 为原点的图坐标。若 \(W(t)\) 是过 \(W\) 的光滑曲线，对 \(w\in W\) 取光滑向量曲线 \(w(t)\in W(t)\)、\(w(0)=w\)，规定其切向量对应的映射为
	\[
		w\longmapsto w'(0)+W\in V/W.
	\]
	这与向量曲线的选择无关：两种选择之差 \(v(t)\in W(t)\) 满足 \(v(0)=0\)，用 \(W(t)\) 的一组光滑基写为 \(\sum_i b_i(t)e_i(t)\)，有 \(b_i(0)=0\)，故 \(v'(0)=\sum_i b_i'(0)e_i(0)\in W\)。若图坐标为 \(A(t)\)、\(A(0)=0\)，可取 \(w(t)=w+A(t)w\)，所得映射就是 \(w\mapsto A'(0)w+W\)。因此它线性、仅依赖切向量，而且给出同构；模 \(W\) 的描述同时表明不依赖补空间。\(k=0\) 或 \(k=n\) 时只有一个子空间，图空间及切空间均为零，结论仍成立。
\end{proof}

% !TeX root = ../Book1.tex
\begin{figure}[htbp]
	\centering
	\begin{tikzpicture}[font=\small,scale=1.05]
		\draw[->,GeometryGray] (-1,0)--(4.2,0) node[right] {\(E\)};
		\draw[->,GeometryGray] (0,-0.65)--(0,3.25) node[above] {\(Q\)};
		\draw[GeometryBlue,very thick] (-0.7,-0.455)--(4,2.6) node[above] {\(W_A\)};
		\draw[->,thick] (0,0)--(2.8,0) node[midway,below] {\(x\)};
		\draw[->,orange!75!black,thick] (2.8,0)--(2.8,1.82) node[midway,right] {\(Ax\)};
		\draw[->,GeometryBlue,thick] (0,0)--(2.8,1.82);
		\fill[GeometryBlue] (2.8,1.82) circle[radius=2pt];
		\node[above left] at (2.8,1.82) {\(x+Ax\)};
		\node[below left] at (0,0) {\(0\)};
		\node[align=center] at (6.5,1.65) {投影 \(W_A\to E\) 可逆\\[4pt]每个 \(x\) 唯一补上 \(Ax\)};
	\end{tikzpicture}
	\caption[Grassmann 流形的图坐标]{图坐标以线性映射 \(A:E\to Q\) 记录子空间 \(W_A\)。图画的是一维横向与一维竖向情形；投影可逆是使用这张坐标图的条件。}
	\label{g1:fig:grassmann-graph}
\end{figure}


\cref{g1:fig:grassmann-graph}表示图坐标的几何意义。我们固定横向的 \(E\)，让子空间 \(W_A\) 倾斜；坐标 \(A\) 记录每个 \(x\in E\) 需要补上的竖向分量。图中只有一个横向和一个竖向维度，但同一构造适用于任意 \(k,n\)。切向量记录的是斜率映射的变化率，而非对子空间内的每个向量任意指定一个独立移动。

\begin{proposition}[Grassmann 流形的度量]\label{g1:prop:grassmann-invariant-metric}
	设 \(\mathbb F=\mathbb R\) 或 \(\mathbb C\)，\(n\ge0\) 及 \(0\leq k\leq n\) 均为整数，\(V=\mathbb F^n\) 分别带标准 Euclid 内积或标准 Hermitian 内积。在 \(\operatorname{Gr}(k,V)\) 上取\cref{g1:thm:grassmann-homogeneous-model}的自然光滑结构，实、复情形分别考虑 \(\operatorname{O}(n)\)、\(\operatorname{U}(n)\) 的标准作用，并用正交投影将 \(\mathrm{T}_W\operatorname{Gr}(k,V)\simeq\operatorname{Hom}_{\mathbb F}(W,V/W)\) 识别为 \(\operatorname{Hom}_{\mathbb F}(W,W^\perp)\)。则
	\begin{equation}\label{g1:eq:grassmann-hilbert-schmidt}
		g_W(S,T)=\operatorname{Re}\operatorname{tr}(S^*T),
		\qquad S,T\in\operatorname{Hom}_{\mathbb F}(W,W^\perp),
	\end{equation}
	定义光滑不变 Riemann 度量。稳定子 \(\operatorname{diag}(C,D)\) 的切表示为 \(T\mapsto DTC^{-1}\)。在 \(\dim_{\mathbb F}E=k\) 的正交分解 \(V=E\oplus Q\) 所给的图坐标 \(A:E\to Q\) 中，对坐标切向量 \(U_1,U_2:E\to Q\)，该度量为
	\begin{equation}\label{g1:eq:grassmann-chart-metric}
		g_A(U_1,U_2)=\operatorname{Re}\operatorname{tr}
		\bigl((I+A^*A)^{-1}U_1^*(I+AA^*)^{-1}U_2\bigr).
	\end{equation}
\end{proposition}
\begin{proof}
	稳定子将 \(x+Tx\) 送到 \(Cx+DTx\)，用 \(x'=Cx\) 改写，图映射变为 \(DTC^{-1}\)，在 \(T=0\) 处也就是所述切表示。由于 \(C,D\) 正交或酉，
	\[
		\operatorname{tr}\bigl((DSC^{-1})^*(DTC^{-1})\bigr)
		=\operatorname{tr}(CS^*TC^{-1})=\operatorname{tr}(S^*T).
	\]
	在正交单位基中，\(\operatorname{tr}(T^*T)\) 是 \(T\) 的所有矩阵分量的模平方之和，故为正定实内积。这是线性映射空间上的平方和内积。\cref{g1:prop:invariant-tensor-isotropy}把它延拓为光滑不变度量；等距变换把正交补送到正交补，所以延拓后的每点公式仍为\eqref{g1:eq:grassmann-hilbert-schmidt}。

	为计算图坐标表达，设 \(S_0=I+A^*A\)、\(T_0=I+AA^*\)。图映射 \(x\mapsto(x,Ax)\) 将 \(W_A\) 的内积拉回为 \(x^*S_0y\)。坐标曲线 \(A+tU\) 中，图向量 \((x,(A+tU)x)\) 的速度为 \((0,Ux)\)。而
	\[
		W_A^\perp=\{(-A^*b,b):b\in Q\}.
	\]
	若 \((0,Ux)\) 的正交投影为 \((-A^*b,b)\)，与任意 \((-A^*c,c)\) 作内积便得 \(c^*T_0b=c^*Ux\)，故 \(b=T_0^{-1}Ux\)。因此两个坐标速度 \(U_1,U_2\) 在输出空间的内积矩阵为 \(U_1^*T_0^{-1}U_2\)。输入空间的度量矩阵是 \(S_0\)，要取正交单位输入基的平方和，就需取 \(\operatorname{tr}(S_0^{-1}U_1^*T_0^{-1}U_2)\)；这可以把一组满足 \(B^*S_0B=I\) 的输入基代入后，用 \(BB^*=S_0^{-1}\) 直接核对。取实部即得\eqref{g1:eq:grassmann-chart-metric}。两个正定矩阵始终可逆，公式中的各项光滑，亦可直接看出整个图上的光滑性。
\end{proof}

当 \(k=1\) 时，Grassmann 流形就是射影空间，\eqref{g1:eq:grassmann-chart-metric}化为\eqref{g1:eq:projective-invariant-metric}：恒等式 \((I+zz^*)^{-1}=I-zz^*/(1+z^*z)\) 可直接相乘验证。更一般地，\(\operatorname{Gr}(2,\mathbb R^4)\) 的一张图由一个 \(2\times2\) 实矩阵的四个分量参数化，实维数为四；\(\operatorname{Gr}(2,\mathbb C^4)\) 使用四个复参数，实维数为八。这里的四维或八维指的是“子空间组成的参数空间”的维数，不能与每个被参数化子空间本身的维数二混为一谈。

% 纲要条目排为编号小节标题，并进入目录与书签；说明段仍保留于纲要注释中。

% !TeX root = ../../Book1.tex
% 纲要标题：### 17.4 非自由作用的局部模型
% 纲要位置：几何学计划纲要.md，第 798 行。
% 正文状态：初稿；尚未正式审核。

% 纲要内容开始
% 
% * 紧 Lie 群固定点附近的线性化
% * 有限稳定子与局部有限群商图
% * 底层商空间和稳定子数据
% * 固有作用的切片定理
% 
% ---
% 
% 纲要内容结束

\section{非自由作用的局部模型}
\label{g1:sec:1704}

自由性使轨道方向能被完全消去，留下普通流形坐标。若一个点仍被某些群元素固定，横截轨道以后，这些元素还会作用在横截面上；商中保留下来的局部识别正由它们产生。本节先把固定点附近的作用写成线性作用，再证明固有性允许将这种局部描述推广到整条轨道的邻域。有限稳定子给出有限群商图，而正维稳定子也能用同一切片构造处理。

\subsection{紧 Lie 群固定点附近的线性化}
\label{g1:subsec:1704:01}

设有限群 \(H\) 光滑作用于 \(M\)，并固定 \(p\)。每个 \(h\in H\) 在 \(p\) 的微分组成线性表示 \(\rho(h)=\mathrm{d}(a_h)_p\)，其中 \(a_h(x)=hx\)。任意坐标 \(c\) 通常不满足 \(c(hx)=\rho(h)c(x)\)。不过，\(\rho(h)^{-1}c(hx)\) 是一份把坐标先移过去、再按线性作用移回来的修正；对所有 \(h\) 平均，就能让每个群元素享有同样地位。

\begin{proposition}[固定点线性化]\label{g1:prop:fixed-point-linearization}
	设有限维实紧 Lie 群 \(H\) 光滑作用于有限维、无边界、Hausdorff、第二可数实光滑流形 \(M\)，且 \(p\) 被 \(H\) 的每个元素固定。令 \(\rho(h)=\mathrm{d}(a_h)_p\)。则存在 \(H\) 不变的开邻域 \(U\ni p\)、\(\mathrm{T}_pM\) 上的 \(H\) 不变内积及关于此内积的零点开球 \(B\)，使某个微分同胚
	\[
		F:U\longrightarrow B\subset \mathrm{T}_pM
	\]
	满足 \(F(p)=0\)、\(\mathrm{d}F_p=\operatorname{id}_{\mathrm{T}_pM}\) 及 \(F(hx)=\rho(h)F(x)\)。有限群作为零维紧 Lie 群包含在此结论中。
\end{proposition}
\begin{proof}
	先取一张以 \(\mathrm{T}_pM\) 为坐标值空间的坐标图 \(c:U_0\to \mathrm{T}_pM\)，使 \(c(p)=0\)、\(\mathrm{d}c_p=I\)；任意坐标图经一个可逆线性变换即可满足后一个条件。我们希望对同一个 \(x\) 的全部 \(hx\) 求坐标，故将定义域缩到
	\[
		U_1=\{x\in M:Hx\subset U_0\}.
	\]
	它含 \(p\) 且 \(H\) 不变。它也是开集：若 \(Hx\subset U_0\)，作用的连续性给每个 \(h\in H\) 一对邻域 \(V_h\ni h\)、\(W_h\ni x\)，使 \(V_hW_h\subset U_0\)。从 \(\{V_h\}\) 取有限子覆盖，再交相应的 \(W_h\)，所得开邻域 \(W\ni x\) 满足 \(HW\subset U_0\)。有限群时，\(U_1\) 也就是有限个开集 \(h^{-1}U_0\) 的交。

	有限群的平均坐标为
	\[
		F(x)=\frac1{|H|}\sum_{h\in H}\rho(h)^{-1}c(hx).
	\]
	一般紧 Lie 群则使用\cref{g1:thm:compact-group-integration}的归一化左右不变积分，定义
	\[
		F(x)=\int_H\rho(h)^{-1}c(hx)\,\mathrm{d}\mu(h).
	\]
	这里逐坐标积分。被积函数光滑，\cref{g1:lem:compact-parameter-integral}允许逐阶在积分下求导，故 \(F\) 光滑：所需的一致控制来自紧参数集 \(H\) 及充分小坐标邻域的紧闭包。有限群时，这个积分直接化为有限个光滑映射之和。对固定 \(k\in H\)，作右平移换元 \(h'=hk\)，得到
	\[
		F(kx)=\int_H\rho(k)\rho(h')^{-1}c(h'x)\,\mathrm{d}\mu(h')
		=\rho(k)F(x).
	\]
	而 \(F(p)=0\)，每一项在 \(p\) 的微分均为 \(\rho(h)^{-1}I\rho(h)=I\)，故 \(\mathrm{d}F_p=I\)。

	由\cref{g1:thm:inverse-function}，\(F\) 在 \(p\) 的一个开邻域 \(V\subset U_1\) 上为到开集的微分同胚。不能直接假定 \(V\) 不变，但可再次使用上面的紧性论证，取开不变邻域 \(V_1=\{x:Hx\subset V\}\)。因为 \(V_1\subset V\)，限制 \(F|_{V_1}\) 仍为微分同胚，且其像开、含零并在 \(\rho(H)\) 下不变。

	最后，将任意正定内积按 \(H\) 平均，得到
	\[
		\langle u,v\rangle_H=\int_H\langle\rho(h)u,\rho(h)v\rangle_0\,\mathrm{d}\mu(h).
	\]
	其正定性与不变性和\cref{g1:thm:compact-group-average}中的论证相同。在开集 \(F(V_1)\) 中取关于此内积的充分小零点球 \(B\)，令 \(U=F^{-1}(B)\)。\(B\) 与 \(U\) 都不变，所需性质随限制保留。
\end{proof}

这说明在固定点附近，我们不仅能得到作用的一阶线性近似，还能选好坐标，使作用本身就等于这个线性表示。若某个 \(h\) 的微分 \(\rho(h)\) 为恒等，它便在这个共同邻域中逐点固定所有点。反过来，\(\rho(h)\ne I\) 时，总有任意靠近 \(p\) 的点被它移动。后面保留稳定子数据时，这一区别会很重要。

\subsection{有限稳定子与局部有限群商图}
\label{g1:subsec:1704:02}

现在 \(p\) 不必被整个 \(G\) 固定，只被稳定子 \(H=G_p\) 固定。\(H\) 在 \(\mathrm{T}_pM\) 上的微分表示保持轨道切空间 \(\mathrm{T}_p(Gp)\)：群元素把轨道送回自身。若作用固有，\cref{g1:prop:proper-action-orbits}保证 \(H\) 紧，因而平均内积给出一个 \(H\) 不变的正交补 \(N\)，满足
\[
	\mathrm{T}_pM=\mathrm{T}_p(Gp)\oplus N.
\]
\(N\) 记录离开轨道的方向。我们需要把它变成真正的横截面，并排除很远的群元素将这张小横截面送回自身的情形。

\begin{lemma}[横截面与返回]\label{g1:lem:proper-slice-construction}
	设有限维实 Lie 群 \(G\) 光滑且固有地作用于有限维实光滑流形 \(M\)，二者无边界、Hausdorff、第二可数。固定 \(p\in M\)，令 \(H=G_p\)、\(\mathfrak g=\mathrm{T}_eG\)、\(\mathfrak h=\mathrm{T}_eH\)。取 \(H\) 不变补空间 \(\mathrm{T}_pM=\mathrm{T}_p(Gp)\oplus N\)，并取线性补空间 \(\mathfrak g=\mathfrak h\oplus\mathfrak m\)。存在 \(H\) 不变嵌入子流形 \(S\ni p\)，使 \(S\) 与 \(N\) 中的一个不变零点开球 \(B\) 等变微分同胚，且具有以下性质：
	\begin{enumerate}
		\item \(\mathrm{T}_pS=N\)，并存在 \(\mathfrak m\) 中的零点邻域 \(V\)，使 \((X,s)\mapsto\exp_G(X)s\) 是 \(V\times S\) 到 \(M\) 中某个开集的微分同胚。
		\item 若 \(s,s'\in S\) 且 \(gs=s'\)，则 \(g\in H\)。因此 \(G_s=H_s\) 对所有 \(s\in S\) 成立。
		\item \(GS\) 在 \(M\) 中开，包含整条轨道 \(Gp\)。包含 \(S\hookrightarrow GS\) 诱导同胚 \(S/H\simeq GS/G\)。
	\end{enumerate}
\end{lemma}
\begin{proof}
	由\cref{g1:thm:closed-subgroup}及\cref{g1:prop:proper-action-orbits}，\(H\) 是紧嵌入 Lie 子群。将\cref{g1:prop:fixed-point-linearization}应用于 \(H\) 在 \(p\) 附近的作用。在得到的等变坐标中取 \(N\) 的小球，再取其逆像 \(S_0\)。这是一张 \(H\) 不变嵌入横截面，\(\mathrm{T}_pS_0=N\)。

	考虑
	\[
		\Psi:\mathfrak m\times S_0\longrightarrow M,\qquad
		\Psi(X,s)=\exp_G(X)s.
	\]
	它在 \((0,p)\) 的微分为 \((X,v)\mapsto X_M(p)+v\)。由\cref{g1:prop:orbit-map-rank}，\(X\mapsto X_M(p)\) 的核是 \(\mathfrak h\)，像是 \(\mathrm{T}_p(Gp)\)，所以限制到 \(\mathfrak m\) 后是到轨道切空间的线性同构。再用直和分解，\(\mathrm{d}\Psi_{(0,p)}\) 可逆。逆函数定理给出局部微分同胚。将其定义域缩成乘积 \(V\times S_1\)，其中 \(V\subset\mathfrak m\) 开且含零，\(S_1\subset S_0\) 仍由一个不变零点小球给出，便有 \(\Psi|_{V\times S_1}\) 单射且微分处处可逆。

	先控制靠近 \(H\) 的群元素。映射
	\[
		H\times\mathfrak m\longrightarrow G,\qquad (h,X)\longmapsto h\exp_G X
	\]
	在每个 \((h,0)\) 处微分可逆，因为左平移后就是 \(\mathfrak h\oplus\mathfrak m\to\mathfrak g\) 的相加映射。因此 \(H\exp_G(V)\) 包含 \(G\) 中一个开邻域 \(W\supset H\)。若 \(g\in W\)、\(s,s'\in S_1\) 且 \(gs=s'\)，可写 \(g=h\exp_G X\)，其中 \(h\in H\)、\(X\in V\)。由 \(S_1\) 不变，\(h^{-1}s'\in S_1\)，于是
	\[
		\Psi(X,s)=h^{-1}s'=\Psi(0,h^{-1}s').
	\]
	\(\Psi\) 的单射性迫使 \(X=0\)，因而 \(g=h\in H\)。这里 \(S_1\) 的不变性让同一个局部逆同时适用于所有 \(h\in H\)，无需逐个群元素选择横截面。

	接着控制不属于 \(W\) 的群元素。若任意小的 \(H\) 不变球形横截面仍有这样的返回，就能取 \(s_j,s'_j\in S_1\)，使 \(s_j,s'_j\to p\)，并取 \(g_j\notin W\) 满足 \(g_js_j=s'_j\)。由\cref{g1:prop:proper-action-criteria}的序列判据，\(g_j\) 有收敛子列，设其极限为 \(g\)。作用连续性给出 \(gp=p\)，即 \(g\in H\subset W\)。因为 \(W\) 开，该子列最终位于 \(W\)，矛盾。因此可将 \(S_1\) 缩成同类不变横截面 \(S\)，使一切满足 \(gs=s'\)、\(s,s'\in S\) 的 \(g\) 都属于 \(W\)，再由上一段得到 \(g\in H\)。这就是所需的返回结论，取 \(s'=s\) 便得 \(G_s=H_s\)。

	\(\Psi(V\times S)\) 开且包含 \(S\)，又包含于 \(GS\)，故其所有 \(G\) 平移的并恰为 \(GS\)，从而 \(GS\) 开。每条其中的轨道都与 \(S\) 相交，返回结论又说明两点 \(s,s'\in S\) 在同一 \(G\) 轨道上，当且仅当它们在同一 \(H\) 轨道上。因此得到连续双射 \(S/H\to GS/G\)。为证明逆连续，任取 \(S/H\) 中开集，设其在 \(S\) 中的原像为开集 \(A\)。\(\Psi(V\times A)\) 在 \(M\) 中开，故其所有 \(G\) 平移的并 \(GA\) 开。\(GA\) 恰是对应商集在 \(GS\) 中的原像，因而该商集在 \(GS/G\) 中开。双射于是为同胚。
\end{proof}

固有性在证明中承担了一个具体任务：如果横截面上的点与其像都逼近 \(p\)，相关群元素不能逃离每个紧集。没有这一步，逆函数定理只能阻止靠近单位元的返回，不能阻止来自群的远处的返回。

\begin{definition}[局部有限群商图]\label{g1:def:orbifold-chart-local}
	设 \(X\) 为拓扑空间，\(U\subset X\) 开。一个\term[orbifoldtu]{局部有限群商图}[orbifold chart]是三元组 \((B,H,\chi)\)：\(B\) 为有限维实内积空间 \(N\) 中的零点开球，有限群 \(H\) 线性地保持 \(B\)，而 \(\chi:B/H\to U\) 是同胚。本节允许表示 \(H\to\operatorname{GL}(N)\) 有非平凡核，也允许其中有反射；表示核非平凡的图称为非有效图，表示忠实的图称为有效图。
\end{definition}

图中的点 \([v]\) 还带有稳定子 \(H_v=\{h\in H:hv=v\}\)；更换同一轨道的代表只把这个子群变为共轭子群，所以其同构类型不变。这正是普通商坐标之外需要保留的数据。一般 局部有限群商图册还要规定不同图之间的相容映射；本章只建立来自给定群作用的局部模型，不展开抽象图册的全部理论。

\begin{corollary}[有限稳定子商的局部模型]\label{g1:cor:finite-stabilizer-local-quotient}
	设有限维实 Lie 群 \(G\) 光滑且固有地作用于有限维无边界实光滑流形 \(M\)，二者 Hausdorff、第二可数，且所有稳定子有限。则每个商点 \([p]\in M/G\) 都有上述 局部有限群商图，其群为 \(G_p\)，其线性作用为 \(G_p\) 在横向空间 \(N\simeq \mathrm{T}_pM/\mathrm{T}_p(Gp)\) 上的表示。局部图的维数为 \(\dim M-\dim G\)，并且图中每点的稳定子与原作用在相应点的稳定子同构。
\end{corollary}
\begin{proof}
	对 \(p\) 应用\cref{g1:lem:proper-slice-construction}。\(S\) 等变微分同胚于 \(N\) 中的球 \(B\)，而 \(S/H\simeq GS/G\)，后者在 \(M/G\) 中开，因此给出所需图。有限群 \(H\) 的 Lie 代数为零，故轨道维数为 \(\dim G\)，横向维数为 \(\dim M-\dim G\)。最后，切片内 \(G_s=H_s\)；在 \(gs\) 处稳定子为 \(gH_sg^{-1}\)，所以图保存了原稳定子的同构类型。
\end{proof}

有限稳定子本身不够。例如 \(\mathbb R\) 沿无理斜率作用于 \(T^2\) 时稳定子平凡，轨道却稠密；商甚至不是 Hausdorff 空间。每个有限群的 Euclid 球商都是 Hausdorff 的，故这样的稠密轨道商不能由上述图描述。固有性与稳定子的有限性分别控制返回现象及残余群的大小，二者不能互相替代。

\subsection{底层商空间和稳定子数据}
\label{g1:subsec:1704:03}

同一个拓扑空间可以有不同的局部群数据。最简单的例子已经出现在二维：一个锥点在底层拓扑上未必是特殊点，但围绕它的有限旋转仍是商结构的一部分。

\begin{example}[有限旋转商]\label{g1:exm:rotation-orbifold-chart}
	令 \(m\geq2\)，\(\mathbb Z/m\mathbb Z\) 在 \(\mathbb C\simeq\mathbb R^2\) 上按 \(z\mapsto\exp(2\pi\mathrm{i}/m)z\) 作用。映射 \(z\mapsto z^m\) 在每条轨道上常值，而且两个点具有相同的 \(m\) 次幂，当且仅当它们在同一轨道上。因此它诱导连续双射
	\[
		\mathbb C/(\mathbb Z/m\mathbb Z)\longrightarrow\mathbb C,\qquad [z]\longmapsto z^m.
	\]
	这是同胚。具体地，若闭集 \(A\subset\mathbb C\) 的像中有收敛序列 \(z_j^m\to w\)，则 \(|z_j|=|z_j^m|^{1/m}\) 有界，存在收敛子列，其极限 \(z\in A\) 满足 \(z^m=w\)。所以幂映射是闭映射，由商拓扑得到上述双射也闭，因而逆连续。

	底空间虽然是一张平面，原点稳定子却是整个 \(\mathbb Z/m\mathbb Z\)，其余点的稳定子平凡。不同的 \(m\)，以及没有群作用的普通平面，具有相同底拓扑而不同稳定子。甚至在坐标 \(w=z^m\) 中，光滑不变函数 \(|z|^2\) 对应 \(|w|^{2/m}\)，后者在原点不光滑。这说明这张拓扑坐标不能直接代替有限群商的光滑数据。
\end{example}

\begin{example}[反射商]\label{g1:exm:reflection-orbifold-chart}
	令二阶群按 \((x,y)\mapsto(x,-y)\) 作用于 \(\mathbb R^2\)。映射 \((x,y)\mapsto(x,|y|)\) 将商同胚地识别为上半平面：每条轨道有唯一的上半平面代表，取该代表的映射连续；商映射限制在上半平面上给出逆映射。边界 \(y=0\) 上的稳定子为整个二阶群，内部点的稳定子平凡。因此，即使被作用的流形无边界，允许反射的有限群商也可能有底拓扑的边界。

	本节的 局部有限群商图允许这种反射。若采用排除反射的文献约定，就必须另加条件，不能将两种定义下的“无边界”直接等同。\cref{g1:fig:finite-quotient-data}将有限旋转商与反射商并列，图中标注的是稳定子，而非额外删去的点。
\end{example}

\begin{figure}[htbp]
	\centering
	% !TeX root = ../Book1.tex
\begin{tikzpicture}[scale=.94, every node/.style={font=\small}]
	\begin{scope}[shift={(-4.4,0)}]
		\fill[blue!9] (0,0) -- (0:1.17) arc (0:120:1.17) -- cycle;
		\draw (0,0) circle (1.17);
		\draw[blue!55] (0,0) -- (0:1.17) (0,0) -- (120:1.17) (0,0) -- (240:1.17);
		\foreach \a in {30,150,270} {\fill[blue!70!black] (\a:.75) circle (1.8pt);}
		\fill (0,0) circle (1.8pt);
		\node[below left] at (0,0) {\(0\)};
		\node at (0,-1.55) {圆盘中的旋转轨道};
	\end{scope}
	\draw[->,thick] (-2.95,0) -- (-1.45,0) node[midway,above] {\(z\mapsto z^m\)};
	\begin{scope}
		\fill[blue!5] (0,0) circle (1.17);
		\draw (0,0) circle (1.17);
		\fill (0,0) circle (1.8pt);
		\node[below] at (0,-.05) {\(\mathbb Z/m\mathbb Z\)};
		\fill (.55,.58) circle (1.8pt);
		\node[above left] at (.55,.58) {\(\{e\}\)};
		\node at (0,-1.55) {旋转商：仍是圆盘};
	\end{scope}
	\begin{scope}[shift={(4.3,-.35)}]
		\fill[blue!5] (-1.17,0) arc (180:0:1.17) -- cycle;
		\draw (-1.17,0) arc (180:0:1.17) -- cycle;
		\draw[blue!65,thick] (-1.17,0) -- (1.17,0);
		\fill (0,0) circle (1.8pt);
		\node[below] at (0,0) {\(\mathbb Z/2\mathbb Z\)};
		\fill (.2,.65) circle (1.8pt);
		\node[above left] at (.2,.65) {\(\{e\}\)};
		\node at (0,-1.2) {反射商：半圆盘};
	\end{scope}
\end{tikzpicture}

	\caption[有限群商中的稳定子]{有限旋转将若干点识别成一个点，商仍同胚于圆盘，但中心保留 \(\mathbb Z/m\mathbb Z\)；反射商是半圆盘，直径上的点保留二阶稳定子。左图以 \(m=3\) 示意。}
	\label{g1:fig:finite-quotient-data}
\end{figure}

还有一种不同的信息损失：某些群元素在整张图中什么也不移动。若 \(K=\ker(H\to\operatorname{GL}(N))\)，则 \(K\) 正规，且 \(H/K\) 在 \(B\) 上有效作用。两者轨道完全相同，因而 \(B/H=B/(H/K)\) 作为拓扑空间自然一致；但是
\[
	(H/K)_v=H_v/K.
\]
将 \(H\) 换为 \(H/K\) 称为这张图的有效化，它同时除掉了每一点稳定子中的 \(K\)，并不保存原来的稳定子数据。

\begin{example}[相同轨道与不同稳定子]\label{g1:exm:ineffective-local-quotient}
	在前面的旋转例子中改用群 \(H=(\mathbb Z/m\mathbb Z)\times(\mathbb Z/k\mathbb Z)\)，其中 \(k\geq2\)，让第二因子恒等作用。商空间及每条轨道均未改变，但原点稳定子变为 \(H\)，非零点稳定子变为 \(\{0\}\times\mathbb Z/k\mathbb Z\)。有效化回到原来的 \(m\) 阶旋转图，恰好删去这些共同的稳定子。这也是本节不自动要求局部群表示忠实的原因。
\end{example}

\subsection{固有作用的切片定理}
\label{g1:subsec:1704:04}

前面的横截面构造并未用到 \(H\) 有限，只用到稳定子紧及作用固有。有限性只是把 \(S/H\) 限制为有限群商。对一般紧稳定子，我们还可以说明整个轨道邻域怎样由一张横截面组装起来：每个 \(g\) 带来一份 \(gS\)，而 \((g,s)\) 与 \((gh^{-1},hs)\) 表示同一个点。这一识别由一个自由作用产生，因此可以使用已经证明的商流形定理。

\begin{theorem}[固有作用的切片定理]\label{g1:thm:proper-action-slice}
	设有限维实 Lie 群 \(G\) 光滑且固有地作用于有限维无边界实光滑流形 \(M\)，二者 Hausdorff、第二可数。对任意 \(p\in M\)，令 \(H=G_p\)，令 \(N\subset \mathrm{T}_pM\) 为某个 \(H\) 不变的轨道切空间补空间，并在 \(N\) 上取 \(H\) 不变内积。存在零点开球 \(B\subset N\)、\(H\) 等变嵌入 \(j:B\to M\) 以及 \(G\) 不变开邻域 \(U\supset Gp\)，满足 \(j(0)=p\)、\(\mathrm{d}j_0\) 为 \(N\hookrightarrow \mathrm{T}_pM\) 的包含，且
	\begin{equation}\label{g1:eq:proper-slice-diffeomorphism}
		G\times_H B\longrightarrow U,\qquad [g,v]\longmapsto g j(v)
	\end{equation}
	为 \(G\) 等变微分同胚。这里
	\[
		G\times_H B=(G\times B)/H,\qquad
		h\cdot(g,v)=(gh^{-1},hv),
	\]
	\(G\) 在其上按 \(a[g,v]=[ag,v]\) 作用。特别地，\(U/G\simeq B/H\)，并且
	\[
		G_{g j(v)}=gH_vg^{-1}.
	\]
\end{theorem}
\begin{proof}
	用\cref{g1:lem:proper-slice-construction}选取 \(S=j(B)\)，令 \(U=GS\)。若事先在 \(N\) 上选定了另一个不变内积，只需在已有线性化坐标域中取关于该内积的更小零点球；引理的局部逆及返回性质在缩小后仍成立。

	先核实 \(G\times_H B\) 是已建立的商流形。上式确实定义左 \(H\) 作用，因为连续作用两次给出
	\[
		h_1\cdot\bigl(h_2\cdot(g,v)\bigr)
		=(gh_2^{-1}h_1^{-1},h_1h_2v)
		=(h_1h_2)\cdot(g,v).
	\]
	它自由，因为 \(gh^{-1}=g\) 强制 \(h=e\)。它固有，因为 \(H\) 紧：对紧集 \(K\subset(G\times B)^2\)，作用映射的原像是 \(H\times\operatorname{pr}_2K\) 中的闭集，故紧。\cref{g1:thm:free-proper-quotient}于是赋予商一个无边界、Hausdorff、第二可数的光滑流形结构，商映射为淹没。\(G\) 的左乘作用与 \(H\) 作用交换，故下降为商上的光滑作用；其光滑性可用商映射的局部截面，在 \(G\times(G\times_HB)\) 的局部坐标中检验。

	光滑映射 \((g,v)\mapsto gj(v)\) 在 \(H\) 轨道上常值，由\cref{g1:prop:quotient-smooth-universal-property}得到光滑映射 \(\overline\Phi:G\times_HB\to U\)。它满射。若 \(g j(v)=g'j(v')\)，则 \(g'^{-1}g\) 将 \(S\) 中一点送到另一点，引理的返回性质给出 \(h=g'^{-1}g\in H\)，且 \(v'=hv\)。于是 \((g',v')=(gh^{-1},hv)\)，两者属于同一个 \(H\) 轨道，证明单射。

	还须检验逆光滑。在 \((e,v)\) 处，原映射沿 \(\mathfrak m\) 方向和 \(\mathrm{T}_vB\) 方向的微分，正是引理中 \((X,s)\mapsto\exp_G(X)s\) 的可逆微分与 \(\mathrm{d}j_v\) 的组合，因而满射。左平移将这点传到每个 \((g,v)\)。原映射通过商映射分解，故 \(\mathrm{d}\overline\Phi\) 处处满射；而
	\[
		\dim(G\times_HB)=\dim G+\dim N-\dim H=\dim M,
	\]
	所以它处处可逆。逆函数定理使 \(\overline\Phi\) 为局部微分同胚，双射性再将各处的局部光滑逆拼成全局光滑逆。这就证明\eqref{g1:eq:proper-slice-diffeomorphism}，等变性由公式直接得到。最后，商空间同胚已由引理证明；稳定子公式由 \(G_{j(v)}=H_{j(v)}=H_v\) 及共轭变换得到。
\end{proof}

这样的 \(S=j(B)\) 称为过 \(p\) 的\term[qiepian]{切片}[slice]。其中 \(N\) 可同构地换成商空间 \(\mathrm{T}_pM/\mathrm{T}_p(Gp)\)，后者的 \(H\) 表示不依赖补空间的选择。不变内积只帮助选择具体代表和小球，并未给 \(M\) 预设一个全局 \(G\) 不变度量。定理也不要求 \(G\) 紧；非紧群的固有性已经足以提供紧稳定子并排除远处返回。

\ProseParagraphBreak

当 \(H=\{e\}\) 时，模型退化为 \(G\times B\)，恢复自由固有作用的局部乘积；当 \(H\) 有限时，商局部为上述有限群商图；当 \(H\) 正维时，\(B/H\) 可以呈现不同维数轨道相邻的现象。例如 \(\operatorname{SO}(2)\) 旋转作用于 \(\mathbb R^2\)，原点的切片就是一个圆盘，商为半开半径区间 \([0,\varepsilon)\)，非零点的轨道是一维圆，原点轨道却是零维。它的底空间恰巧也像反射商的一维截面，但原点稳定子是整条圆周群，不能用有限稳定子图保留同样的数据。

\ProseParagraphBreak

本节已证明的是光滑固有作用在一条轨道附近的切片模型。关于不同稳定子类型如何组成分层、不同切片如何组织为完整 局部有限群商图册，以及带边界或其他范畴中的切片版本，还需要相应的定义与证明。后续使用本节结果时，可以直接调用\eqref{g1:eq:proper-slice-diffeomorphism}及其商 \(B/H\)，而不能由此把所有非自由商都当作普通流形。


\begin{chaptersummary}
	沿 \(\exp(tX)\cdot p\) 取初速度得到完备基本场 \(X_M\)，左作用和正号约定使 \([X_M,Y_M]=-[X,Y]_M\)。稳定子 \(G_p\) 闭，其 Lie 代数是轨道微分的核；轨道自然具有 \(G/G_p\) 的浸入流形结构，切空间为 \(\mathfrak g/\mathfrak g_p\)。传递作用使这一模型成为整个流形的微分同胚描述。

	\ProseParagraphBreak

	固有性控制两端留在紧区域内的全部群参数，使稳定子紧、轨道闭嵌入，并排除横截面的远处返回。与自由性合用，得到 \(G\times S\cong GS\) 的局部乘积，因而 \(M/G\) 是唯一使投影成为淹没的光滑商。轨道不变光滑映射可唯一光滑下降；只要求投影光滑不足以保证这一性质。

	\ProseParagraphBreak

	左平移延拓单位元处的正密度，紧性给出有限总量；模函数的紧像只能平凡，于是归一化积分同时左右不变。平均度量保留正定性，紧参数积分保证光滑性。齐性空间上的不变张量则由基点处的稳定子不变张量唯一决定，不必对整个群都作积分。

	\ProseParagraphBreak

	球面和射影空间由“固定向量”与“固定直线”的不同稳定子产生。Grassmann 流形以线性映射的图作局部坐标，\(\mathrm{T}_W\operatorname{Gr}(k,V)\cong\operatorname{Hom}_{\mathbb F}(W,V/W)\)。选定内积后可用 \(W^\perp\) 代表商空间，从而直接计算稳定子作用及平方和度量；复情形的实维数须乘二。

	\ProseParagraphBreak

	紧群固定点附近可通过平均坐标把作用完全线性化。对固有作用，这给出轨道邻域 \(G\times_{G_p}B\)，商邻域为 \(B/G_p\)。有限稳定子产生局部有限群商图，其中反射可能带来底拓扑边界，非有效核则保留额外稳定子。底空间的同胚类型不能代替这些群数据；一般正维稳定子的商也不自动是普通流形或有限群商。
\end{chaptersummary}

\BookExercises


\begin{exercise}\label{g1:ex:conjugation-orbit-dimension}
	令 \(G=\operatorname{GL}(n,\mathbb R)\) 在 \(M_n(\mathbb R)\) 上按共轭作用。求 \(X\in\mathfrak{gl}(n,\mathbb R)\) 对应的基本场、矩阵 \(A\) 的稳定子 Lie 代数及共轭轨道切空间。若 \(A\) 可实对角化，互异特征值的重数为 \(m_1,\ldots,m_s\)，求轨道维数。
\end{exercise}
\begin{solution}
	对 \(\exp(tX)A\exp(-tX)\) 求导，得到
	\[
		X_M(A)=XA-AX.
	\]
	稳定子是可逆中心化子 \(G_A=\{g:gA=Ag\}\)，由\cref{g1:prop:orbit-map-rank}，其 Lie 代数为 \(\mathfrak g_A=\{X:XA=AX\}\)。\cref{g1:prop:orbit-immersion}进一步给出
	\[
		\mathrm{T}_A(G\cdot A)=\{XA-AX:X\in M_n(\mathbb R)\}.
	\]
	这里切空间对应轨道的自然浸入结构，不需要先假定轨道闭。

	取特征空间直和 \(\mathbb R^n=E_1\oplus\cdots\oplus E_s\)，\(A|_{E_i}=\lambda_iI\)。将 \(X\) 写为分块矩阵 \((X_{ij})\)，交换关系在第 \((i,j)\) 块成为 \((\lambda_j-\lambda_i)X_{ij}=0\)。不同特征值使非对角块都为零；对角块 \(X_{ii}\in\operatorname{End}(E_i)\) 可以任取。因此
	\[
		\dim\mathfrak g_A=\sum_{i=1}^s m_i^2,
		\qquad \dim(G\cdot A)=n^2-\sum_{i=1}^s m_i^2.
	\]
	特别地，互异实特征值给出维数 \(n^2-n\)，标量矩阵给出零维单点轨道。后一个例子也说明群维数大，并不意味着它在每个对象上都产生很多运动方向。
\end{solution}

\begin{exercise}\label{g1:ex:affine-action-fundamental-fields}
	令 \(G=(0,\infty)\times\mathbb R\)，乘法为 \((a,b)(c,d)=(ac,b+ad)\)，作用于 \(\mathbb R\) 的公式为 \((a,b)\cdot x=ax+b\)。取单位元处的基 \(H=\partial_a\)、\(E=\partial_b\)。求基本场、各点的稳定子及稳定子 Lie 代数，并把 \(\mathbb R\) 写成 \(G\) 的齐性空间。说明该作用是否自由、是否有效。
\end{exercise}
\begin{solution}
	单参数子群分别为 \(\exp(tH)=(e^t,0)\)、\(\exp(tE)=(1,t)\)，所以
	\[
		H_{\mathbb R}=x\partial_x,\qquad E_{\mathbb R}=\partial_x,
		\qquad [H_{\mathbb R},E_{\mathbb R}]=-\partial_x.
	\]
	前章算得 \([H,E]=E\)，因此这个例子核对了正号基本场约定下的反同态公式。固定 \(x\) 的条件是 \(ax+b=x\)，故
	\[
		G_x=\{(a,(1-a)x):a>0\},\qquad
		\mathfrak g_x=\mathbb R(H-xE).
	\]
	后一个等式也可以从 \((uH+vE)_{\mathbb R}(x)=(ux+v)\partial_x\) 的核直接得到。

	平移可把 \(0\) 送到任意 \(x\)，所以作用传递。令 \(H_0=G_0=\{(a,0):a>0\}\)，则\cref{g1:prop:transitive-action-model}给出 \(G/H_0\cong\mathbb R\)，对应为 \((a,b)H_0\mapsto b\)。也可直接核对：右乘 \((c,0)\) 只改变第一坐标，\(b\) 不变，且 \(x\mapsto(1,x)H_0\) 是光滑逆。每点稳定子都有一维，故作用不自由。如果 \(ax+b=x\) 对所有 \(x\) 成立，先令 \(x=0\) 得 \(b=0\)，再令 \(x=1\) 得 \(a=1\)，所以作用有效。
\end{solution}

\begin{exercise}\label{g1:ex:circle-equivariant-maps}
	设 \(a,b>0\)。加法群 \(\mathbb R\) 按平移作用于 \(\mathbb R/a\mathbb Z\) 和 \(\mathbb R/b\mathbb Z\)。求两者之间的所有 \(\mathbb R\) 等变光滑映射，并求每个映射的纤维基数及微分。它们何时为微分同胚？
\end{exercise}
\begin{solution}
	设 \(F(0+a\mathbb Z)=c+b\mathbb Z\)。等变性强制
	\[
		F(t+a\mathbb Z)=t+c+b\mathbb Z.
	\]
	代表元 \(t\) 与 \(t+a\) 给出同一个目标值，当且仅当 \(a\in b\mathbb Z\)，即 \(a=kb\)，其中 \(k\) 为正整数。这与\cref{g1:prop:equivariant-homogeneous-maps}中的 \(a\mathbb Z\subseteq b\mathbb Z\) 完全一致。若此条件不成立，就没有等变映射；若成立，参数 \(c\) 按模 \(b\mathbb Z\) 取值，给出全部映射。

	固定目标点 \(s+b\mathbb Z\)，其原像由 \(s-c+jb\)、\(j\in\mathbb Z\)，在模 \(a\mathbb Z\) 的意义下给出。因为 \(a=kb\)，恰有 \(k\) 个不同类，可取 \(0\le j<k\)。在局部提升到实直线的坐标中，\(F\) 为 \(t\mapsto t+c\)，故微分是恒等线性映射，处处为局部微分同胚。它为微分同胚恰好当 \(k=1\)，也即 \(a=b\)。注意局部微分可逆并不能排除整体上出现多个原像。
\end{solution}

\begin{exercise}\label{g1:ex:equivariant-automorphisms-normalizer}
	设 \(H\) 为 Lie 群 \(G\) 的闭子群，记正规化子 \(N_G(H)=\{a\in G:aHa^{-1}=H\}\)。证明 \(G/H\) 的全部 \(G\) 等变微分同胚与 \(N_G(H)/H\) 一一对应。核对复合次序，并说明如何选取公式使这一对应保持群乘法。
\end{exercise}
\begin{solution}
	由\cref{g1:prop:equivariant-homogeneous-maps}，等变映射必为 \(F_a(gH)=gaH\)。若它还是双射，则对任意 \(k\in G\)，
	\[
		k\cdot F_a(eH)=F_a(eH)
		\quad\Longleftrightarrow\quad F_a(kH)=F_a(eH)
		\quad\Longleftrightarrow\quad k\in H.
	\]
	另一方面，\(aH\) 的稳定子是 \(aHa^{-1}\)，所以必须有 \(aHa^{-1}=H\)。反之，若 \(a\in N_G(H)\)，则 \(F_a\) 与 \(F_{a^{-1}}\) 都有定义且光滑，互为逆。\(F_a=F_b\) 当且仅当 \(aH=bH\)，得到所需集合对应。

	\(N_G(H)\) 是子群，且 \(H\) 在其中正规，故 \(N_G(H)/H\) 是群。按普通函数复合顺序，
	\[
		(F_a\circ F_b)(gH)=gbaH=F_{ba}(gH).
	\]
	因此 \(aH\mapsto F_a\) 反转乘法。改定义 \(R_a(gH)=ga^{-1}H\)，则
	\[
		(R_a\circ R_b)(gH)=gb^{-1}a^{-1}H
		=g(ab)^{-1}H=R_{ab}(gH),
	\]
	从而 \(aH\mapsto R_a\) 是到等变微分同胚群的群同构。公式中逆元的位置由右乘与函数复合的次序决定。
\end{solution}


\begin{exercise}\label{g1:ex:screw-action-quotient}
	令 \(R_t\) 表示平面上转角 \(t\) 的旋转矩阵。加法群 \(\mathbb R\) 在 \(\mathbb R^2\times\mathbb R\) 上作用为
	\[
		t\cdot(v,z)=(R_tv,z+t).
	\]
	证明此作用自由且固有，求其基本向量场，并用一个明确的光滑不变量把商识别为 \(\mathbb R^2\)。
\end{exercise}
\begin{solution}
	旋转矩阵满足 \(R_{s+t}=R_sR_t\)，因此给出的公式确为光滑作用。若 \(t\cdot(v,z)=(v,z)\)，第三坐标给出 \(t=0\)，故自由。若 \((v_j,z_j)\to(v,z)\)，且 \(t_j\cdot(v_j,z_j)\to(w,u)\)，则由第三坐标 \(z_j+t_j\to u\)，得到 \(t_j\to u-z\)。因此\cref{g1:prop:proper-action-criteria}的序列条件成立，作用固有。

	以 \(1\in\mathbb R\) 为 Lie 代数生成元，在 \(t=0\) 求导，得到基本场
	\[
		X=-y\partial_x+x\partial_y+\partial_z.
	\]
	考虑 \(F(v,z)=R_{-z}v\)。由于 \(R_{-(z+t)}R_t=R_{-z}\)，它在轨道上常值。每条轨道恰与平面 \(z=0\) 相交一次，交点为 \((F(v,z),0)\)。因此下降映射 \(\bar F:(\mathbb R^2\times\mathbb R)/\mathbb R\to\mathbb R^2\) 双射且光滑，其逆 \(w\mapsto[(w,0)]\) 也是光滑映射。更直接地，
	\[
		(t,w)\longmapsto(R_tw,t),\qquad
		(v,z)\longmapsto(z,R_{-z}v)
	\]
	是互逆的全局光滑乘积坐标。轨道的螺旋形状在这些坐标中变为固定 \(w\)、只改变 \(t\)。
\end{solution}

\begin{exercise}\label{g1:ex:diagonal-proper-action}
	设有限维实 Lie 群 \(G\) 自由且固有地光滑作用于无边界光滑流形 \(M\)，并任意光滑作用于无边界光滑流形 \(N\)；所有流形均为有限维、Hausdorff、第二可数。证明对角作用
	\[
		g\cdot(p,n)=(gp,gn)
	\]
	在 \(M\times N\) 上自由且固有。特别地，取 \(M=G\) 且作用为左平移，证明 \((G\times N)/G\) 微分同胚于 \(N\)。
\end{exercise}
\begin{solution}
	若 \(g\) 固定 \((p,n)\)，它固定 \(p\)，所以 \(g=e\)。若 \((p_j,n_j)\) 收敛且 \((g_jp_j,g_jn_j)\) 收敛，则 \(p_j\) 与 \(g_jp_j\) 都收敛。原作用的固有性给出 \(g_j\) 的收敛子列，因此对角作用也满足序列判据。

	左平移在 \(G\) 上自由且固有：如果 \(a_j\to a\)、\(g_ja_j\to b\)，就有 \(g_j=(g_ja_j)a_j^{-1}\to ba^{-1}\)。在 \(G\times N\) 上令
	\[
		F(a,n)=a^{-1}n.
	\]
	则 \(F(ga,gn)=a^{-1}n\)，故它下降为光滑映射 \(\bar F:(G\times N)/G\to N\)。每条轨道包含唯一一点 \((e,a^{-1}n)\)，所以 \(\bar F\) 为双射，且逆映射 \(n\mapsto[(e,n)]\) 光滑。这证明所需微分同胚。注意不需要 \(G\) 在 \(N\) 上的作用自由或固有；第一因子的平移已控制了群参数。
\end{solution}

\begin{exercise}\label{g1:ex:closed-subgroup-action-quotients}
	设有限维实 Lie 群 \(G\) 自由且固有地光滑作用于有限维、Hausdorff、第二可数的无边界光滑流形 \(M\)，\(H\subseteq G\) 为闭子群。证明限制的 \(H\) 作用也自由且固有，并证明
	\[
		F:M/H\longrightarrow M/G,\qquad F(Hp)=Gp
	\]
	是光滑满射淹没。说明若把 \(H\) 换成带自身 Lie 群拓扑的非闭浸入子群，固有性不再能由同一论证保证，并给出反例。
\end{exercise}
\begin{solution}
	闭子群定理使 \(H\) 成为具有子空间拓扑的嵌入 Lie 子群。自由性直接继承。记 \(\Theta_G(g,p)=(gp,p)\)，则对每个紧集 \(C\subseteq M\times M\)，有
	\[
		\Theta_H^{-1}(C)=\Theta_G^{-1}(C)\cap(H\times M).
	\]
	右侧是紧集的闭子集，因而紧。这里 \(H\) 的闭性与子空间拓扑同时被使用。

	记商投影为 \(\pi_H\)、\(\pi_G\)。映射 \(\pi_G\) 光滑且在 \(H\) 轨道上常值，所以光滑下降性质给出唯一光滑 \(F\)，满足 \(F\circ\pi_H=\pi_G\)。由 \(\pi_G\) 满射可知 \(F\) 满射。对 \(p\in M\) 求微分，得到
	\[
		\mathrm{d}F_{Hp}\circ\mathrm{d}(\pi_H)_p=\mathrm{d}(\pi_G)_p.
	\]
	右侧满射，因此 \(\mathrm{d}F_{Hp}\) 满射，证明 \(F\) 是淹没。定义 \(F\) 无须假定 \(H\) 正规：这里比较的是轨道空间，尚未给 \(M/H\) 定义 \(G/H\) 的群作用。

	反例取 \(G=M=T^2\)，作用为左平移；它自由且因 \(G\) 紧而固有。取\cref{g1:exm:irrational-torus-subgroup}的无理斜率浸入子群 \(H_\alpha\cong\mathbb R\)。限制作用正是本章的无理流。其底空间 \(M\) 紧，但搬运集 \(T(M,M)=H_\alpha\cong\mathbb R\) 非紧，故作用不固有。
\end{solution}

\begin{exercise}\label{g1:ex:smooth-quotient-without-submersion}
	令 \(\mathbb R\) 以 \(t\cdot(x,y)=(x+t,y)\) 作用于 \(\mathbb R^2\)，轨道集记为 \(Q\)。通常用 \(y\) 给 \(Q\) 坐标。现改用
	\[
		u:Q\longrightarrow\mathbb R,\qquad u([x,y])=y^3
	\]
	作为全局坐标，赋予 \(Q\) 一份光滑结构。证明这份结构与原结构具有相同拓扑，商投影仍光滑，但不是淹没，并举出一个不再能光滑下降的不变光滑函数。这说明商流形定理的唯一性为什么要求投影为淹没。
\end{exercise}
\begin{solution}
	通常的商坐标 \([x,y]\mapsto y\) 是到 \(\mathbb R\) 的同胚；\(s\mapsto s^3\) 又是 \(\mathbb R\) 的同胚，逆为连续函数 \(u\mapsto\sqrt[3]{u}\)。因此 \(u\) 仍是商拓扑的一张全局坐标，作为单张图可以规定光滑结构。

	在新坐标下，投影 \(\pi:\mathbb R^2\to Q\) 的表达为 \(u\circ\pi(x,y)=y^3\)，故光滑。但其微分为 \(3y^2\,\mathrm{d}y\)，在直线 \(y=0\) 上为零，所以不是到一维流形的淹没。原来的函数 \(f(x,y)=y\) 光滑且平移不变，其唯一集合意义的下降在新坐标下为 \(\bar f(u)=\sqrt[3]{u}\)，在零点不可微。因此它不能光滑下降。两份结构之间，原商到新商的恒等映射表达为 \(s\mapsto s^3\)，反向为立方根，故恒等双射不是微分同胚。
\end{solution}

\begin{exercise}\label{g1:ex:closed-orbits-nonproper-action}
	加法群 \(\mathbb R\) 在 \(M=\mathbb R^2\setminus\{(0,0)\}\) 上作用为
	\[
		t\cdot(x,y)=(e^t x,e^{-t}y).
	\]
	证明作用自由，且每条轨道在 \(M\) 中闭；再证明作用不固有，轨道商也不是 Hausdorff 空间。由此说明“每条轨道闭”不能替代固有性。
\end{exercise}
\begin{solution}
	因为 \((x,y)\ne(0,0)\)，至少有一个坐标非零。若 \(e^t x=x\)、\(e^{-t}y=y\)，由该非零坐标得到 \(t=0\)，所以自由。

	乘积 \(xy\) 与每个非零坐标的符号在作用下不变。如果 \(xy=c\ne0\)，固定 \(x\) 的符号以后，曲线 \(xy=c\) 的相应分支就是一条轨道，因为两个同号的 \(x\) 值之比可写成 \(e^t\)。这条分支在 \(M\) 中闭：任何有限极限仍满足 \(xy=c\ne0\)，且 \(x\) 不会从正变负或从负变正。若 \(xy=0\)，轨道是四条坐标半轴之一；它们在去掉原点的 \(M\) 中也闭。因此所有轨道都闭。

	取 \(p_j=(e^{-j},1)\)、\(t_j=j\)，则
	\[
		p_j\longrightarrow p=(0,1),\qquad
		t_jp_j=(1,e^{-j})\longrightarrow q=(1,0).
	\]
	群参数 \(t_j\) 没有收敛子列，故作用不固有。另一方面，\(p,q\) 位于不同轨道。若商中分别包含 \([p]\)、\([q]\) 的开邻域为 \(U,V\)，其原像是分别包含 \(p,q\) 的开集。充分大的 \(j\) 满足 \(p_j\in\pi^{-1}(U)\)、\(t_jp_j\in\pi^{-1}(V)\)，而这两点属于同一轨道，故 \(U\cap V\ne\varnothing\)。所以这两个商点不能由不交开邻域分开。
\end{solution}


\begin{exercise}\label{g1:ex:finite-nonorthogonal-average}
	令 \(C=\left(\begin{smallmatrix}0&-1\\1&-1\end{smallmatrix}\right)\)。验证 \(C^3=I\)，并把 \(\mathbb R^2\) 上的标准内积对群 \(\{I,C,C^2\}\) 平均。求所得矩阵 \(Q\)，直接验证 \(C^{\mathsf T}QC=Q\)，并说明为什么 \(C\) 不保持标准内积，却能保持一个正定内积。
\end{exercise}
\begin{solution}
	相乘得到
	\[
		C^2=\begin{pmatrix}-1&1\\-1&0\end{pmatrix},\qquad C^3=I.
	\]
	有限群的归一化积分为各元素取权重 \(1/3\)，故
	\[
		Q=\frac13\bigl(I+C^{\mathsf T}C+(C^2)^{\mathsf T}C^2\bigr)
		=\frac13\begin{pmatrix}4&-2\\-2&4\end{pmatrix}.
	\]
	它的特征值为 \(2/3\)、\(2\)，所以正定。直接相乘有
	\[
		QC=\frac13\begin{pmatrix}-2&-2\\4&-2\end{pmatrix},\qquad
		C^{\mathsf T}QC=\frac13\begin{pmatrix}4&-2\\-2&4\end{pmatrix}=Q.
	\]
	另一方面，\(C^{\mathsf T}C=\left(\begin{smallmatrix}1&-1\\-1&2\end{smallmatrix}\right)\ne I\)，故标准内积不被保持。正交性依赖选定内积；有限群平均构造了一个与该群作用相容的新内积。它的单位圆在原坐标中是一条椭圆，\(C\) 保持的是这条椭圆。
\end{solution}

\begin{exercise}\label{g1:ex:average-projection-fixed-vectors}
	设 \(K\) 为紧 Lie 群，\(\rho:K\to\operatorname{GL}(V)\) 为有限维实光滑表示。逐矩阵分量定义 \(P=\int_K\rho(k)\,\mathrm{d}\mu_K(k)\)。证明 \(P^2=P\)，且 \(\operatorname{im}P=V^K\)，其中 \(V^K=\{v:\rho(k)v=v\text{ 对所有 }k\in K\}\)。进一步证明
	\[
		\dim V^K=\int_K\operatorname{tr}\rho(k)\,\mathrm{d}\mu_K(k),
	\]
	并求圆周群在 \(\mathbb R^2\) 上的标准旋转表示的 \(P\)。
\end{exercise}
\begin{solution}
	矩阵分量是连续函数，紧性保证所有积分有限。对 \(h\in K\)，利用左、右不变积分分别得到
	\[
		\rho(h)P=\int_K\rho(hk)\,\mathrm{d}\mu_K(k)=P,
		\qquad P\rho(h)=P.
	\]
	所以 \(Pv\) 被每个 \(\rho(h)\) 固定，\(\operatorname{im}P\subseteq V^K\)。若 \(v\in V^K\)，则 \(Pv=\int_Kv\,\mathrm{d}\mu_K=v\)，反向包含也成立，而且 \(P^2v=Pv\)。这证明 \(P\) 是到固定向量空间的投影，并非任意向量的平均都非零。

	投影给出直和 \(V=\operatorname{im}P\oplus\ker P\)，因为 \(v=Pv+(v-Pv)\)，两子空间相交为零。在适应该直和的基中，\(P\) 的矩阵为 \(\operatorname{diag}(I,0)\)，故 \(\operatorname{tr}P=\dim V^K\)。迹是对角分量的有限和，可以与积分交换，得到题中公式。旋转表示的矩阵为 \(R_t=\left(\begin{smallmatrix}\cos t&-\sin t\\\sin t&\cos t\end{smallmatrix}\right)\)，各分量在一周上的平均均为零，所以 \(P=0\)，也与没有非零向量被全部旋转固定一致。
\end{solution}

\begin{exercise}\label{g1:ex:sphere-invariant-tensors}
	对 \(n\geq1\)，证明 \(S^n\) 上每个 \(\operatorname{SO}(n+1)\) 不变 Riemann 度量都是单位球面标准度量的正常数倍。再求所有 \(\operatorname{SO}(n+1)\) 不变向量场，并说明 \(n=1\) 与 \(n\geq2\) 的区别。
\end{exercise}
\begin{solution}
	由\cref{g1:prop:sphere-homogeneous-model,g1:prop:invariant-tensor-isotropy}，问题等价于求 \(\mathbb R^n\) 上的 \(\operatorname{SO}(n)\) 不变正定内积及固定向量。若 \(n\geq2\)，\(\operatorname{SO}(n)\) 在单位球面 \(S^{n-1}\) 上传递：把任一单位向量补成正向正交单位基，再将两组这样的基对应，即得所需矩阵。因而一个不变内积 \(b\) 的二次型在标准单位向量上取同一值 \(c>0\)，按齐次性有 \(b(v,v)=c|v|^2\)。极化公式
	\[
		2b(v,w)=b(v+w,v+w)-b(v,v)-b(w,w)
	\]
	于是给出 \(b(v,w)=c\langle v,w\rangle\)。当 \(n=1\) 时，一维空间上的任意正定内积本来就是该形式。基点数据的唯一延拓证明了全部 \(n\geq1\) 的度量结论。

	若 \(n\geq2\)，任意非零 \(v\in\mathbb R^n\) 都能与一个垂直向量张成二维平面；在该平面转动 \(\pi\)，在正交补上取恒等，得到 \(\operatorname{SO}(n)\) 中把 \(v\) 变为 \(-v\) 的元素。因此固定向量只有零，不变向量场也只有零。若 \(n=1\)，稳定子 \(\operatorname{SO}(1)\) 平凡，基点切空间的每个向量都能延拓。将 \(S^1\) 写为 \(p=(\cos t,\sin t)\)，所有不变向量场正是
	\[
		X_p=c(-\sin t,\cos t),\qquad c\in\mathbb R.
	\]
	右侧也可写成 \(cJp\)，所以与角度分支的选择无关。
\end{solution}

\begin{exercise}\label{g1:ex:projective-line-radius}
	给 \(\mathbb{CP}^1\) 配备\eqref{g1:eq:projective-invariant-metric}的度量。证明仿射坐标 \(z=x+iy\) 中
	\[
		g=\frac{\mathrm{d}x^2+\mathrm{d}y^2}{(1+x^2+y^2)^2},
	\]
	并证明映射
	\[
		F([1:z])=\frac1{1+|z|^2}
		\left(x,y,\frac{|z|^2-1}{2}\right)
	\]
	延拓为 \(\mathbb{CP}^1\) 到半径 \(1/2\) 的 Euclid 球面的等距微分同胚。
\end{exercise}
\begin{solution}
	一维复坐标中，\((1+|z|^2)\overline uv-(\overline uz)(\overline zv)=\overline uv\)，故\eqref{g1:eq:projective-invariant-metric}立即给出题中的实度量。记 \(s=x^2+y^2\)、\(r=1+s\)。直接计算
	\[
		|F|^2=\frac{s+(s-1)^2/4}{(1+s)^2}=\frac14,
	\]
	所以像落在指定球面。有限 \(z\) 的像不含北极 \((0,0,1/2)\)，其余点 \((X,Y,Z)\) 的逆为 \(z=(X+iY)/(1/2-Z)\)。这证明该坐标片到去掉北极的球面的双射。

	在另一张射影图中令 \(w=1/z=a+ib\)。对 \(w\ne0\)，同一公式变为
	\[
		F([w:1])=\frac1{1+|w|^2}
		\left(a,-b,\frac{1-|w|^2}{2}\right),
	\]
	右侧在 \(w=0\) 光滑，取值恰为北极。由两张图及球面的相应坐标逆，可得整体微分同胚。

	最后核对度量。对 \(\xi=(u,v)\in\mathbb R^2\) 令 \(q=xu+yv\)。\(F\) 的前两个分量的微分为 \(\xi/r-2q(x,y)/r^2\)，第三个分量的微分为 \(2q/r^2\)。故
	\begin{align*}
		|\mathrm{d}F(\xi)|^2
		&=\frac{|\xi|^2}{r^2}-\frac{4q^2}{r^3}
		+\frac{4q^2s}{r^4}+\frac{4q^2}{r^4}\\
		&=\frac{|\xi|^2}{r^2}.
	\end{align*}
	极化后得到两切向量的内积也一致；北极处由另一张光滑图中的连续性成立。因此它是等距微分同胚。球半径为 \(1/2\)，具体反映本节射影度量的归一化，不能在保持同一个度量的同时把它写成单位球面。
\end{solution}

\begin{exercise}\label{g1:ex:grassmann-complement-isometry}
	设 \(V=\mathbb F^n\) 带标准内积，\(\mathbb F=\mathbb R\) 或 \(\mathbb C\)。证明正交补映射
	\[
		c:\operatorname{Gr}(k,V)\longrightarrow\operatorname{Gr}(n-k,V),
		\qquad W\longmapsto W^\perp
	\]
	是光滑微分同胚，且对\eqref{g1:eq:grassmann-hilbert-schmidt}定义的两边度量是等距映射。明确写出 \(\mathrm{d}c_W\) 在切空间 \(\operatorname{Hom}_{\mathbb F}(W,W^\perp)\) 上的表达。
\end{exercise}
\begin{solution}
	固定正交分解 \(V=W\oplus W^\perp\)，邻近的子空间写成 \(W_A=\{x+Ax:x\in W\}\)。向量 \(u+v\)、\(u\in W\)、\(v\in W^\perp\) 垂直于全部 \(x+Ax\)，当且仅当 \(u=-A^*v\)，故
	\[
		W_A^\perp=\{-A^*v+v:v\in W^\perp\}.
	\]
	在目标使用交换后的分解 \(V=W^\perp\oplus W\)，正交补映射的坐标表达就是 \(A\mapsto-A^*\)。它在实情形线性，在复情形实线性，因而都是实光滑映射；不能把复共轭误认为复线性。又因 \((W^\perp)^\perp=W\)，其逆也是同一正交补构造，故为微分同胚。

	在 \(A=0\) 处求导得到
	\[
		\mathrm{d}c_W(T)=-T^*:
		W^\perp\longrightarrow W.
	\]
	因此对 \(S,T:W\to W^\perp\)，目标度量给出
	\[
		\operatorname{Re}\operatorname{tr}\bigl((-S^*)^*(-T^*)\bigr)
		=\operatorname{Re}\operatorname{tr}(ST^*)
		=\operatorname{Re}\operatorname{tr}(S^*T).
	\]
	最后一步可在正交单位基中展开为 \(\sum_{i,j}\operatorname{Re}(s_{ij}\overline{t_{ij}})\)，与 \(\sum_{i,j}\operatorname{Re}(\overline{s_{ij}}t_{ij})\) 相等；这也适用于两个迹分别在不同维数空间上计算的情形。于是每点的内积都被保持，得到等距性。\(k=0\) 或 \(k=n\) 时两空间均为单点，结论按零切空间解释成立。
\end{solution}


\begin{exercise}\label{g1:ex:compact-action-fixed-set}
	设紧 Lie 群 \(K\) 光滑作用于无边界、Hausdorff、第二可数光滑流形 \(M\)，定义共同固定点集
	\[
		M^K=\{p\in M:kp=p\text{ 对所有 }k\in K\}.
	\]
	证明 \(M^K\) 闭，并且每个固定点附近的固定点集都是嵌入子流形，切空间为
	\[
		\mathrm{T}_p(M^K)=\{v\in \mathrm{T}_pM:\mathrm{d}(a_k)_pv=v\text{ 对所有 }k\in K\}.
	\]
	说明各连通分支的维数可以不同，并用 \(\mathbb{RP}^2\) 上的反射给出一个例子。
\end{exercise}
\begin{solution}
	对每个 \(k\)，集合 \(\{p:kp=p\}\) 是映射 \(p\mapsto(kp,p)\) 对闭对角线的原像，故闭；这些闭集的交 \(M^K\) 也闭。固定 \(p\in M^K\)，由\cref{g1:prop:fixed-point-linearization}，在 \(p\) 的某个不变邻域中有等变坐标 \(F:U\to B\subset \mathrm{T}_pM\)，且 \(\mathrm{d}F_p=I\)。令 \(V\) 为题目右侧的固定向量子空间。等变性与 \(F\) 的单射性给出
	\[
		F(U\cap M^K)=B\cap V.
	\]
	右侧是线性子空间中的开球，因而是一张嵌入坐标切片；其零点切空间是 \(V\)，经 \(\mathrm{d}F_p=I\) 送回即得所求公式。固定点集的维数因此在每个这样的局部片内恒定，从而在每个连通分支上恒定，但不同分支未必有相同维数。

	具体地，在 \(\mathbb{RP}^2\) 上令二阶群的生成元作用为
	\[
		[x:y:z]\longmapsto[x:y:-z].
	\]
	这是由线性反射诱导的光滑作用。一个射影点固定，当且仅当存在非零实数 \(\lambda\)，使 \((x,y,-z)=\lambda(x,y,z)\)。若 \((x,y)\ne(0,0)\)，必有 \(\lambda=1\)、\(z=0\)，得到射影直线 \(\{z=0\}\simeq\mathbb{RP}^1\)。若 \((x,y)=(0,0)\)，得到单点 \([0:0:1]\)。二者不相交，固定点集因此有一个一维连通分支和一个零维连通分支。故未经补充不能把整个固定点集说成同一固定维数的子流形。
\end{solution}

\begin{exercise}\label{g1:ex:weighted-circle-local-quotients}
	设正整数 \(a,b\geq1\)，令 \(S^1\) 光滑作用于 \(S^3\subset\mathbb C^2\)：
	\[
		\lambda\cdot(z_1,z_2)=(\lambda^a z_1,\lambda^b z_2).
	\]
	求两坐标均非零处及两条坐标圆上的稳定子。计算 \(p=(1,0)\)、\(q=(0,1)\) 处的横向稳定子表示，说明这些局部图何时有效；再求整个作用的核以及有效化后两处的旋转阶数。
\end{exercise}
\begin{solution}
	作用保持 \(|z_1|^2+|z_2|^2\)，并且是光滑作用。由于 \(S^1\) 紧，作用固有，因此可用本节的切片结果。记 \(d=\gcd(a,b)\)，并以 \(\mu_r=\{\lambda\in S^1:\lambda^r=1\}\) 表示 \(r\) 次单位根群。如果 \(z_1z_2\ne0\)，稳定条件为 \(\lambda^a=\lambda^b=1\)，恰好等价于 \(\lambda^d=1\)：一向由 \(d\mid a,b\)，另一向由整数 \(u,v\) 使 \(ua+vb=d\)。故一般点稳定子为 \(\mu_d\)。当 \(z_2=0\) 时稳定子是 \(\mu_a\)，当 \(z_1=0\) 时稳定子是 \(\mu_b\)。

	在 \(p=(1,0)\)，切空间为
	\[
		\mathrm{T}_pS^3=\{(\mathrm{i}t,w):t\in\mathbb R,\ w\in\mathbb C\}.
	\]
	轨道方向由 \((\mathrm{i}a,0)\) 张成，故可取横向空间 \(N_p=\{(0,w):w\in\mathbb C\}\)。\(\lambda\in\mu_a\) 在这个复直线上的作用为 \(w\mapsto\lambda^b w\)。还可以直接取小圆盘上的横截面
	\[
		j_p(w)=\bigl(\sqrt{1-|w|^2},w\bigr),
	\]
	它对 \(\mu_a\) 等变；这也在实际坐标中显示了上述表示。其核为 \(\mu_a\cap\mu_b=\mu_d\)，像是阶数 \(a/d\) 的旋转群。交换两坐标，在 \(q\) 得到 \(\mu_b\) 的表示 \(w\mapsto\lambda^a w\)，核仍为 \(\mu_d\)，像的阶数为 \(b/d\)。两张图都有效当且仅当 \(d=1\)。

	一个 \(\lambda\) 若在整个 \(S^3\) 上恒等作用，检验 \(p,q\) 即得 \(\lambda^a=\lambda^b=1\)，所以全局作用核也是 \(\mu_d\)。有效化后，一般点稳定子变平凡，两处特殊图的旋转阶数分别为 \(a/d\)、\(b/d\)。阶数为 \(1\) 时该处就是普通圆盘图。有效化前则不能删除共同的 \(\mu_d\)，否则会改变每个点的稳定子。
\end{solution}

\begin{exercise}\label{g1:ex:finite-action-faithful-linearization}
	设有限群 \(K\) 有效地光滑作用于连通无边界光滑流形 \(M\)。证明对任意 \(p\in M\)，稳定子的表示 \(K_p\to\operatorname{GL}(\mathrm{T}_pM)\) 忠实。这里“有效”指只有单位群元素在整个 \(M\) 上恒等作用。若去掉连通性，给出反例。
\end{exercise}
\begin{solution}
	设 \(h\in K_p\)，且 \(\mathrm{d}(a_h)_p=I\)。对有限循环群 \(\langle h\rangle\) 应用固定点线性化，可知 \(h\) 在 \(p\) 的某个邻域上恒等。我们还要把这条局部结论传播到整个连通流形；一般光滑映射在开集上恒等并不足够，\(h\) 的有限阶性质在下一步仍要使用。

	令
	\[
		A=\{x\in M:hx=x,\ \mathrm{d}(a_h)_x=I\}.
	\]
	这个集合非空。它闭：若 \(x_j\in A\) 趋于 \(x\)，连续性先给 \(hx=x\)；在 \(x\) 附近的一张共同坐标图中，\(a_h\) 的导数矩阵连续，所以其在 \(x\) 的值也是单位矩阵。它也开：对每个 \(x\in A\)，再次对 \(\langle h\rangle\) 在固定点 \(x\) 应用线性化，得到 \(h\) 在 \(x\) 的整个小邻域上恒等，因而该邻域均在 \(A\) 中。连通性迫使 \(A=M\)，所以 \(h\) 恒等作用。全局有效性给出 \(h=e\)，证明忠实性。

	若 \(M=\mathbb R\sqcup\mathbb R\)，让二阶群在第一份直线上恒等作用，在第二份上按 \(t\mapsto-t\) 作用。这是有效光滑作用，但第一份上每个点都被整个群固定，且稳定子的微分表示全为恒等。反例表明，连通性负责把局部不动传播到全局，不能省略。
\end{solution}

\begin{exercise}\label{g1:ex:effective-transitive-ineffective-slice}
	令 \(G=\operatorname{SO}(3)\)，\(r=\operatorname{diag}(-1,-1,1)\)，\(H=\{I,r\}\)。考虑 \(G\) 在齐性空间 \(M=G/H\) 上的左作用。证明这个作用有效、固有，所有稳定子均为二阶群，但每个点的横向切片都是零维且带平凡的稳定子表示。说明它为什么不与上一题矛盾，又为什么不能将该商只按一个普通点理解而保留所有稳定子信息。
\end{exercise}
\begin{solution}
	\(H\) 是有限闭子群，因此\cref{g1:thm:homogeneous-space}给出 \(G/H\) 的光滑流形结构，维数为 \(3\)。左作用传递；在 \(gH\) 处，稳定子由 \(agH=gH\) 算得为 \(gHg^{-1}\)，所以均为二阶群。\(G\) 紧，故作用固有。

	整个作用的核是 \(\bigcap_{g\in G}gHg^{-1}\)，特别地包含于 \(H\)。非单位元素 \(r\) 是绕第三坐标轴的半周旋转；取把第三坐标轴转到第一坐标轴的 \(g\)，则 \(grg^{-1}=\operatorname{diag}(1,-1,-1)\ne r\)。两个二阶子群 \(H\) 与 \(gHg^{-1}\) 的交只有单位元，故全局核平凡，作用有效。

	传递性使每个点的轨道就是整个 \(M\)，所以 \(\mathrm{T}_p(Gp)=\mathrm{T}_pM\)，横向空间 \(N\) 为零空间。其零点球只有一个点，整个二阶稳定子都在这一点上平凡作用；故切片图非有效。上一题讨论的是有限群作用在整个连通流形上的微分表示，这里作用群 \(G\) 正维，而且我们只保留商掉轨道方向以后的横向表示，假设与表示对象都不同。

	底层商 \(M/G\) 是一个点，而切片图还记录 \(H\simeq\mathbb Z/2\mathbb Z\)。把该图有效化后成为带平凡群的一个点，底拓扑不变，稳定子却丢失。这一例子说明全局作用有效，并不保证横向切片表示有效。
\end{solution}

